\documentclass[10pt,a4paper]{article}
\usepackage[margin=2.2cm]{geometry}
\usepackage{fontspec}
\setmonofont{DejaVu Sans Mono}[Scale=0.84]
\newfontfamily\unifont{DejaVu Serif}[Scale=0.9]
\usepackage{amsmath,amssymb,mathtools}
\setcounter{MaxMatrixCols}{20}
\usepackage{longtable,booktabs,array}
\usepackage{fancyvrb}
\usepackage[dvipsnames]{xcolor}
\usepackage[colorlinks=true,linkcolor=NavyBlue,urlcolor=NavyBlue]{hyperref}
\usepackage{microtype}
\usepackage{needspace}
\usepackage{tocloft}\setlength{\cftbeforesecskip}{0.45em}
\setlength{\parindent}{0pt}\setlength{\parskip}{0.55em}
\setlength{\emergencystretch}{8em}
\renewcommand{\arraystretch}{1.15}
\allowdisplaybreaks
\begin{document}
\title{7+1: the octonions and the Clifford algebra of 7+1 spacetime\\[0.4em]\large A real realization of the octonions, and a real 16×16 realization of the Clifford algebra of 7+1 spacetime, built from — and mirroring — the split octonions and the 16×16 Dirac matrices of the notebook 4+4.nb}
\author{Repository: github.com/once-ere/7plus1 — prepared for Nash\\ Notebook: 7+1.nb (Wolfram Language 15.0.1), evaluated Tue 27 Oct 2026 19:31:35\\ Document generated by tools/make\_docs.py on 2026-10-27}
\date{}
\maketitle
\begin{abstract}
Starting from the real 8×8 matrices $\tau_A$ and the real 16×16 Dirac matrices $T_A$ of 4+4.nb (split octonions and the Clifford algebra $Cl(4,4)$), this document and the notebook 7+1.nb construct (i) a real realization of the octonions whose multiplication constants come from exactly the same vielbein formula as the split-octonion constants of 4+4.nb; (ii) real 16×16 generators of the Clifford algebra of 7+1 spacetime (one time, seven space dimensions), built only from the matrices of 4+4.nb; (iii) a proof that in the sign convention of 4+4.nb no real 16×16 realization exists, together with the minimal complexification and a real 32×32 realization; (iv) examples that test everything; (v) in sections 7–10 of the notebook, the Dirac equation of the 16-component quantized fermion field dirac16complex with these real Dirac matrices, its chiral form as two coupled 8-component equations, and the 16-component and chiral Kohn–Sham theories of that field coupled to itself (section 11); (vi) in sections 11–19 of the notebook, gravity on the 7+1 spacetime — curvature, the Lovelock tensors (computed with xAct exactly as in the reference notebook Einstein-Lovelock-7+1-Nash.nb, and independently), generalized frame matrices, the canonical spin connection that couples gravity to dirac16complex, and the primordial, modern and far-future epochs — the nine-channel self-coupling of dirac16complex with its Kohn–Sham theory in gravitational fields, and the Standard Model on the world brane with the portals of dirac16complex and neutrino flavour conversion (section 12); (vii) in sections 20–23 of the notebook, the coupled field equations of the whole model with every interaction written as an equation, and accurate approximate solutions of them in the three gravity regimes — primordial (with its cubic, Gauss–Bonnet and Einstein sub-regimes), modern and far future — each compared with a high-accuracy numerical solution (section 13). All 399 checks of the notebook are True.
\end{abstract}
\tableofcontents
\Needspace{12\baselineskip}\section*{1. Results at a glance}\addcontentsline{toc}{section}{1. Results at a glance}
\textbf{Octonions from the split-octonion matrices.} With $\sigma=\tau_1\tau_2\tau_3=\tau_4\tau_5\tau_6\tau_7$ (the 8×8 matrix \texttt{σ} of 4+4.nb),
\begin{equation*}
\tau^{O}_0=1_8,\qquad \tau^{O}_A=\tau_A\ \ (A=1,2,3),\qquad \tau^{O}_A=\tau_A\,\sigma\ \ (A=4,5,6,7),\qquad \bar\tau^{O}_A=(\tau^{O}_A)^{\mathsf T}.
\end{equation*}
These are real, $\tau^{O}_A\bar\tau^{O}_B+\tau^{O}_B\bar\tau^{O}_A=2\,\delta_{AB}\,1_8$, and the vielbein formula of 4+4.nb turns them into the multiplication table of the octonions $\mathbb{O}$ (section 6).

\textbf{Real Clifford algebra of 7+1 spacetime.} With the 16×16 Dirac matrices $T_A$ of 4+4.nb and $S_{44}=T_4T_5T_6T_7$,
\begin{equation*}
\Gamma^{(1,7)}_0=T_0,\qquad \Gamma^{(1,7)}_i=T_0\,T_i\ \ (i=1,2,3),\qquad \Gamma^{(1,7)}_k=T_0\,T_k\,S_{44}\ \ (k=4,5,6,7),
\end{equation*}
\begin{equation*}
\Gamma^{(1,7)}_A\Gamma^{(1,7)}_B+\Gamma^{(1,7)}_B\Gamma^{(1,7)}_A=2\,\eta^{(1,7)}_{AB}\,1_{16},\qquad \eta^{(1,7)}=\mathrm{diag}(+1,-1,-1,-1,-1,-1,-1,-1).
\end{equation*}
All entries are $0,\pm1$. Time ($x_0$, the octonion unit) squares to $+1$, the seven space directions to $-1$ — exactly like the octonion basis: $\eta^{(1,7)}$ is simply the squares of the octonion units, $e_0^2=+1$, $e_i^2=-1$, and $\eta^{(1,7)}_{AB}=\mathrm{Re}(e_Ae_B)$. The quadratic form is the real part of the octonion square, $\mathrm{Re}(x\,x)=x_0^2-x_1^2-\dots-x_7^2$; the algebra is $Cl(1,7)\cong\mathbb{R}(16)$ (section 7.2). The spacetime has one time-like and seven space-like dimensions; the tensor package xAct (section 12 of 7+1.nb) declares it, as the author’s notebook Einstein-Lovelock-7+1-Nash.nb does, with the signature list $\{7,1,0\}$: seven space-like dimensions, one time-like dimension, none degenerate (section 2).

\textbf{The sign convention of 4+4.nb.} 4+4.nb writes metrics as "space-like $=+1$, time-like $=-1$". In that convention the 7+1 metric is $\eta^{(7,1)}=\mathrm{diag}(-1,+1,\dots,+1)=-\eta^{(1,7)}$ and the relation $\Gamma_A\Gamma_B+\Gamma_B\Gamma_A=2\eta^{(7,1)}_{AB}1_{16}$ defines $Cl(7,1)\cong\mathbb{H}(8)$, which has NO real 16×16 realization, whatever the order or signs of the indices (theorem of section 7.4). The minimal complexification multiplies only the time block by $i$ and keeps the exact block form of 4+4.nb:
\begin{align*}
\Gamma^{(7,1)}_A&=\begin{pmatrix}0&\bar\tau^{C}_A\\ \tau^{C}_A&0\end{pmatrix},\qquad \tau^{C}_0=\bar\tau^{C}_0=i\,1_8,\quad \tau^{C}_k=\tau^{O}_k,\ \bar\tau^{C}_k=\bar\tau^{O}_k\ (k=1,\dots,7), \\
\Gamma^{(7,1)}_A\Gamma^{(7,1)}_B+\Gamma^{(7,1)}_B\Gamma^{(7,1)}_A&=2\,\eta^{(7,1)}_{AB}\,1_{16}.
\end{align*}
A real realization of $Cl(7,1)$ needs 32×32 matrices: $\widehat\Gamma^{(7,1)}_A=\varepsilon\otimes\Gamma^{(1,7)}_A$ with $\varepsilon=\left(\begin{smallmatrix}0&1\\-1&0\end{smallmatrix}\right)$ (section 7.5).

\textbf{The Dirac equation of dirac16complex and its Kohn–Sham theories} (sections 7–10 of the notebook; section 11 below). With the real $\Gamma^{(1,7)}_A$ and $\Gamma^{A}=\eta_{(1,7)}^{AA}\Gamma^{(1,7)}_A$, the 16-component field $\Psi$ of complex anticommuting operators obeys $(i\Gamma^{A}\partial_A-m)\Psi=0$; its chiral 8-component parts $\Psi_L=(\Psi_\uparrow+i\Psi_\downarrow)/\sqrt2$, $\Psi_R=(i\Psi_\uparrow+\Psi_\downarrow)/\sqrt2$ obey
\begin{equation*}
\big(i\,\partial_0+\tau^{O}_i\,\partial_i\big)\Psi_L=m\,\Psi_R,\qquad \big(i\,\partial_0-\tau^{O}_i\,\partial_i\big)\Psi_R=m\,\Psi_L,
\end{equation*}
coupled by the mass. The field coupled to itself by scalar, pseudoscalar, vector and axial-vector contact interactions has a 16-component Kohn–Sham theory and a chiral (8+8-component) Kohn–Sham theory with a complex Kohn–Sham mass.

\textbf{Gravity, the self-coupled field in gravitational fields, and the Standard Model} (sections 11–19 of the notebook; section 12 below). On a metric $g_{\mu\nu}(X)$ of the 7+1 spacetime — one time-like and seven space-like dimensions, declared in xAct, as in the reference notebook Einstein-Lovelock-7+1-Nash.nb, with the signature list $\{7,1,0\}$ (seven space-like, one time-like, none degenerate), and components with the signs of $\eta^{(1,7)}=\mathrm{diag}(+1,-1,\dots,-1)$, simply the squares of the octonion basis ($g_{00}>0$) — the Lovelock tensors of order $k$, with the generalized Kronecker delta $\delta^{\mu_1\cdots\mu_n}_{\nu_1\cdots\nu_n}=\det[\delta^{\mu_i}_{\nu_j}]$ and the Riemann tensor $R^{\gamma\delta}{}_{\alpha\beta}$, are
\begin{equation*}
E^{(k)\mu}{}_\nu=-\frac{1}{2^{k+1}}\,\delta^{\mu\alpha_1\beta_1\cdots\alpha_k\beta_k}_{\nu\gamma_1\delta_1\cdots\gamma_k\delta_k}\,R^{\gamma_1\delta_1}{}_{\alpha_1\beta_1}\cdots R^{\gamma_k\delta_k}{}_{\alpha_k\beta_k}\qquad(E^{(1)}=G,\ E^{(4)}\equiv0\ \text{in eight dimensions}),
\end{equation*}
and the generalized frame matrices $\Gamma_\mu(X)=e^A{}_\mu(X)\Gamma^{(1,7)}_A$ with the canonical spin connection $\Omega_\mu$ have a vanishing total covariant derivative, so that dirac16complex couples to gravity through
\begin{equation*}
\nabla_\mu\Gamma_\nu=\partial_\mu\Gamma_\nu-\{{}^{\lambda}_{\mu\nu}\}\Gamma_\lambda+[\Omega_\mu,\Gamma_\nu]=0,\qquad \big(i\,\Gamma^\mu(X)(\partial_\mu+\Omega_\mu)-m\big)\Psi=0 .
\end{equation*}
The field coupled to itself through all nine ranks of the Clifford algebra has a Kohn–Sham theory in gravitational fields, solved in the three cosmological epochs with its back-reaction on the expansion; the Standard Model lives on the brane $x_4=\dots=x_7=0$ and meets dirac16complex through portals. Sections 11–19 contain 148 checks of their own.

\textbf{The coupled field equations and their solutions in the three gravity regimes} (sections 20–23 of the notebook; section 13 below). Every interaction of the model — the four Lovelock terms, dirac16complex with its mass, nine self-couplings and spin connection, a stabilizing potential of the auxiliary space, the brane tension, the Standard Model on the brane and the portals — is a term of one action, whose variation gives coupled field equations; for the homogeneous metric $ds^2=dx_0^2-a^2d\vec x^2-b^2d\vec y^2$ they become
\begin{equation*}
\sum_{k=0}^{3}\alpha_kE^{(k)0}{}_0=\kappa\varepsilon,\qquad \sum_{k=0}^{3}\alpha_kE^{(k)1}{}_1=-\kappa p_a,\qquad \sum_{k=0}^{3}\alpha_kE^{(k)4}{}_4=-\kappa p_b .
\end{equation*}
In the primordial regime the vacuum solution is exact (an algebraic curve and one quadrature), with Kasner limits $(p_1,q_1)=(0.54692,-0.16019)$ (Einstein) and $(p_2,q_2)=(1.5611,-0.42086)$ (Gauss–Bonnet) whose closed-form corrections are accurate to 0.00032 and 0.00085; a dirac16complex gas makes the universe matter-dominated near $x_0=0$ and, without a stabilizing potential, makes the auxiliary space re-expand. With the potential the radion is caught and carries 0.991 of the energy (the moduli problem of the minimal model); the modern brane cosmology has the composite analytic solution accurate to 0.00022 (age 13.791 Gyr), and the far future is de Sitter with a frozen auxiliary space. Sections 20–23 contain 48 checks of their own.

\textbf{Verification.} The notebook 7+1.nb contains 399 checks; evaluated in 15.0.1 for Microsoft Windows (64-bit) (July 2, 2026) it returns True for every one of them, with 0 message cells, in 271.569 s. Independent Python programs reproduce every matrix, table and number that the notebook prints (section 9).

\Needspace{12\baselineskip}\section*{2. Coordinates, metrics and conventions}\addcontentsline{toc}{section}{2. Coordinates, metrics and conventions}
4+4.nb and 7+1.nb use the same eight coordinate symbols $X=\lbrace x_0,x_1,x_2,x_3,x_4,x_5,x_6,x_7\rbrace$ (Mathematica: \texttt{X = \{x0, x1, x2, x3, x4, x5, x6, x7\}}; never $x_1,\dots,x_8$). Their meaning differs:

{\small
\begin{longtable}{>{\raggedright\arraybackslash}p{\dimexpr 0.1702\linewidth-2\tabcolsep\relax}>{\raggedright\arraybackslash}p{\dimexpr 0.4043\linewidth-2\tabcolsep\relax}>{\raggedright\arraybackslash}p{\dimexpr 0.4255\linewidth-2\tabcolsep\relax}}
\toprule
\textbf{coordinate} & \textbf{meaning in 4+4.nb (signature 4+4)} & \textbf{meaning in 7+1.nb (signature 7+1)} \\
\midrule\endfirsthead
\toprule
\textbf{coordinate} & \textbf{meaning in 4+4.nb (signature 4+4)} & \textbf{meaning in 7+1.nb (signature 7+1)} \\
\midrule\endhead
\endfoot
\bottomrule\endlastfoot
$x_0$ & hidden space-like dimension & ordinary time; octonion unit $e_0=1$ \\
$x_1,x_2,x_3$ & ordinary 3-space & ordinary 3-space; imaginary units $e_1,e_2,e_3$ \\
$x_4$ & ordinary time & auxiliary space-like; $e_4=\ell$ (Cayley–Dickson doubling unit) \\
$x_5,x_6,x_7$ & three auxiliary time-like dimensions ("superluminal deflating time") & auxiliary space-like; $e_5,e_6,e_7$ \\*
\end{longtable}}
Metrics (Mathematica names in brackets):
\begin{align*}
\eta^{(4,4)}&=\mathrm{diag}(+1,+1,+1,+1,-1,-1,-1,-1) && [\texttt{η4488}]\ \text{(4+4.nb: space-like }+1\text{, time-like }-1) \\
\eta^{(7,1)}&=\mathrm{diag}(-1,+1,+1,+1,+1,+1,+1,+1) && [\texttt{η71}]\ \text{(7+1 in the sign convention of 4+4.nb)} \\
\eta^{(1,7)}&=-\eta^{(7,1)}=\mathrm{diag}(+1,-1,\dots,-1) && [\texttt{η17}]\ \text{(7+1: the squares of the octonion basis)} \\
g^{O}&=1_8 && [\texttt{gO}]\ \text{(norm of the octonions)}
\end{align*}
\textbf{The signature: $\{7,1,0\}$ and $\mathrm{diag}(+1,-1,\dots,-1)$.} The 7+1 spacetime has one time-like dimension ($x_0$, the direction of the octonion unit $e_0=1$) and seven space-like dimensions ($x_1,x_2,x_3$ ordinary space, $x_4,\dots,x_7$ auxiliary space). Two different objects describe this, and neither may be replaced by the other. (i) The tensor package xAct, used in section 12 of 7+1.nb exactly as in the author’s notebook Einstein-Lovelock-7+1-Nash.nb, declares the metric with \texttt{DefMetric[\{7, 1, 0\}, g71[-a, -b], CD, …]}: the list $\{7,1,0\}$ declares seven space-like dimensions, one time-like dimension and no degenerate dimension — the numbers of entries $+1$ and $-1$ of $\eta^{(7,1)}=\mathrm{diag}(-1,+1,\dots,+1)$ (\texttt{η71}), the same spacetime written in the sign convention of 4+4.nb (space-like $=+1$, time-like $=-1$). (ii) The matrix $\eta^{(1,7)}=-\eta^{(7,1)}=\mathrm{diag}(+1,-1,\dots,-1)$ (\texttt{η17}) is simply the squares of the octonion basis: $e_0^2=+1$ for the unit $e_0=1$ (time) and $e_i^2=-1$ for the seven imaginary units $e_1,\dots,e_7$ (space); more generally it is the real part of the product of two basis units, and it gives the quadratic form of the octonions:
\begin{gather*}
\eta^{(1,7)}_{AB}=\mathrm{Re}(e_Ae_B),\qquad e_0^2=+1,\qquad e_i^2=-1\ \ (i=1,\dots,7),\\ \mathrm{Re}(x\,x)=x_0^2-x_1^2-\cdots-x_7^2=\eta^{(1,7)}_{AB}\,x^Ax^B\qquad (x=x^Ae_A\in\mathbb O).
\end{gather*}
$\eta^{(1,7)}$ gives the signs of the Clifford relation $\Gamma^{(1,7)}_A\Gamma^{(1,7)}_B+\Gamma^{(1,7)}_B\Gamma^{(1,7)}_A=2\,\eta^{(1,7)}_{AB}\,1_{16}$ of the real 16×16 Dirac matrices: time squares to $+1_{16}$ and each of the seven space directions to $-1_{16}$, exactly like the octonion units (in the notation $Cl(p,q)$, with $p$ generators squaring to $+1$ and $q$ to $-1$, the algebra $Cl(1,7)\cong\mathbb R(16)$). It also gives the signs of the metric components: $g_{00}>0$ along the time-like direction, negative squares along the space-like directions. Every component computed by 7+1.nb follows from the component matrix of the metric; the declared list fixes for xAct the sign of the determinant of the metric, \texttt{SignDetOfMetric[g71]} $=-1$, which is the same in both sign conventions, since $\det \eta^{(1,7)}=\det \eta^{(7,1)}=-1$. 7+1.nb checks \texttt{SignatureOfMetric[g71]} $=\{7,1,0\}$, \texttt{SignDetOfMetric[g71]} $=-1$, $\mathrm{Re}(e_Ae_B)=\eta^{(1,7)}_{AB}$ and $\mathrm{Re}(x\,x)=\eta^{(1,7)}_{AB}x^Ax^B$.

Conventions. Indices $A,B,K,L\in\lbrace0,\dots,7\rbrace$ label coordinates, octonion units and Clifford generators; $i,j\in\lbrace1,\dots,7\rbrace$ label imaginary units and space directions; $h\in\lbrace1,2,3\rbrace$. $1_n$ is the identity matrix, $M^{\mathsf T}$ the transpose, $\delta_{AB}$ the Kronecker delta, $\lbrace X,Y\rbrace=XY+YX$, and $\left(\begin{smallmatrix}a&b\\c&d\end{smallmatrix}\right)$ a matrix of blocks. The index 8 on a Clifford generator denotes the ordered product of all eight, as in 4+4.nb ($T_8=T_0T_1\cdots T_7$). $Cl(p,q)$ always means: $p$ generators square to $+1$ and $q$ to $-1$, i.e. $v^2=Q(v)$ with $Q$ of signature $(p,q)$ (for the real 7+1 generators the signs of the squares of the octonion basis, hence $Cl(1,7)$; this label is not the xAct list $\{7,1,0\}$, which counts seven space-like and one time-like dimension). Mathematica lists are 1-based, so the metric entry $\eta_{AB}$ is written \texttt{η[[A+1, B+1]]} in the notebook.

\Needspace{12\baselineskip}\section*{3. Notation used in this document}\addcontentsline{toc}{section}{3. Notation used in this document}
Every symbol of this document, with its Mathematica name in 7+1.nb and its definition. Names in the first group are reproduced exactly from 4+4.nb (where some of them carried 2-D names made with the Notation package: $\bar\tau_A$ was a symbolized tau-bar, $T_A$ was \texttt{T16\textasciicircum{}A[A]}, $E_{44}$ was \texttt{E\textasciicircum{}A\_a[h]}); the plain one-line names avoid confusing them with powers or indices.

{\small
\begin{longtable}{>{\raggedright\arraybackslash}p{\dimexpr 0.1809\linewidth-2\tabcolsep\relax}>{\raggedright\arraybackslash}p{\dimexpr 0.2128\linewidth-2\tabcolsep\relax}>{\raggedright\arraybackslash}p{\dimexpr 0.6064\linewidth-2\tabcolsep\relax}}
\toprule
\textbf{symbol} & \textbf{Mathematica} & \textbf{definition / meaning} \\
\midrule\endfirsthead
\toprule
\textbf{symbol} & \textbf{Mathematica} & \textbf{definition / meaning} \\
\midrule\endhead
\endfoot
\bottomrule\endlastfoot
$\sigma$ & \texttt{σ} & $\left(\begin{smallmatrix}0&1_4\\1_4&0\end{smallmatrix}\right)$, split metric on spinor space (4+4.nb \texttt{σ88}) \\
$s_h,\ t_h$ & \texttt{s4by4[h]}, \texttt{t4by4[h]} & self-dual / anti-self-dual 4×4 matrices, $(s_h)_{pq}=\epsilon_{hpq4}-(\delta_{p4}\delta_{qh}-\delta_{ph}\delta_{q4})$, $(t_h)_{pq}=\epsilon_{hpq4}+(\delta_{p4}\delta_{qh}-\delta_{ph}\delta_{q4})$ \\
$\zeta_A$ & \texttt{ζ[A]} & $\zeta_0=1_8$, $\zeta_h=\left(\begin{smallmatrix}0&s_h\\s_h&0\end{smallmatrix}\right)$, $\zeta_{7-h}=\left(\begin{smallmatrix}0&t_h\\-t_h&0\end{smallmatrix}\right)$, $\zeta_7=\zeta_1\cdots\zeta_6$ \\
$\tau_A,\ \bar\tau_A$ & \texttt{τ[A]}, \texttt{τbar[A]} & $\tau_A=\zeta_A$, $\bar\tau_A=\sigma\tau_A^{\mathsf T}\sigma$ (split octonions) \\
$T_A$ & \texttt{T16[A]} & $\left(\begin{smallmatrix}0&\bar\tau_A\\\tau_A&0\end{smallmatrix}\right)$, Dirac matrices of 4+4.nb; $T_8=T_0\cdots T_7$ \\
$P_L,\ P_R$ & \texttt{PL44}, \texttt{PR44} & $(1_{16}\mp T_8)/2$ \\
$u$ & \texttt{unit} & $(1,0,0,0,1,0,0,0)/\sqrt2$, the unit spinor (\texttt{uBasis[[8]]}, \texttt{hUSE = 8}) \\
$E_{44}$ & \texttt{E44[8]} & split vielbein, rows $u^{\mathsf T}\sigma\bar\tau_B$ \\
$\varepsilon_B$ & \texttt{eA[B]} & split-octonion units, $B=1,\dots,8$ \\
$\tau^{O}_A,\ \bar\tau^{O}_A$ & \texttt{τO[A]}, \texttt{τObar[A]} & octonion matrices, section 6.1 \\
$E^{O}$ & \texttt{EO} & octonion vielbein, rows $u^{\mathsf T}\bar\tau^{O}_B$ \\
$L^{O}_A$ & \texttt{LO[A]} & $E^{O}\tau^{O}_A(E^{O})^{\mathsf T}$, left multiplication by $e_A$ \\
$e_A$ & \texttt{e[A]} & octonion units, $A=0,\dots,7$, $e_0=1$ \\
$N(x),\ \bar x,\ \mathrm{Re}\,x$ & \texttt{oNorm2}, \texttt{oConj}, \texttt{oRe} & norm, conjugate, real part \\
$T^{O}_A$ & \texttt{T16O[A]} & $\left(\begin{smallmatrix}0&\bar\tau^{O}_A\\\tau^{O}_A&0\end{smallmatrix}\right)$, generators of $Cl(8,0)$ \\
$S_{44}$ & \texttt{S44} & $T_4T_5T_6T_7=\mathrm{diag}(\sigma,\sigma)$ \\
$\Gamma^{(1,7)}_A$ & \texttt{Γ17[A]} & real generators of the 7+1 Clifford algebra, metric $\eta^{(1,7)}$ \\
$\Gamma^{(7,1)}_A$ & \texttt{Γ71[A]} & complex generators, metric $\eta^{(7,1)}$ \\
$\widehat\Gamma^{(7,1)}_A$ & \texttt{Γ71R32[A]} & real 32×32 generators, metric $\eta^{(7,1)}$ \\
$\Gamma^{(5,3)}_A,\ \Gamma^{(0,8)}_A$ & \texttt{Γ53[A]}, \texttt{Γ08[A]} & real generators of $Cl(5,3)$ and $Cl(0,8)$ (section 8.6) \\
$\Sigma_{AB}$ & \texttt{Σ17[A,B]} & $(\Gamma^{(1,7)}_A\Gamma^{(1,7)}_B-\Gamma^{(1,7)}_B\Gamma^{(1,7)}_A)/4$ \\
$\varepsilon$ & \texttt{ε2} & $\left(\begin{smallmatrix}0&1\\-1&0\end{smallmatrix}\right)$ \\
$\sigma_1,\ \sigma_3$ & — & $\left(\begin{smallmatrix}0&1\\1&0\end{smallmatrix}\right)$, $\left(\begin{smallmatrix}1&0\\0&-1\end{smallmatrix}\right)$ \\
$\Gamma^{(1,7)}_8$ & \texttt{Γ17[8]} & $\Gamma^{(1,7)}_0\Gamma^{(1,7)}_1\cdots\Gamma^{(1,7)}_7=\left(\begin{smallmatrix}0&-1_8\\1_8&0\end{smallmatrix}\right)$, the volume element \\
$Q_i,\ Q_j,\ Q_k$ & \texttt{Qi}, \texttt{Qj}, \texttt{Qk} & $\varepsilon\otimes1_{16}$, $\sigma_1\otimes\Gamma^{(1,7)}_8$, $\sigma_3\otimes\Gamma^{(1,7)}_8$ (quaternion units, section 7.5) \\
$\chi^{(7,1)},\ \chi^{(1,7)}$ & \texttt{χ71}, \texttt{χ17} & $-i\,\Gamma^{(7,1)}_8$ and $i\,\Gamma^{(1,7)}_8$, chirality matrices (section 7.6) \\
$\vec x$ & \texttt{xImag} & imaginary part $x_1e_1+\dots+x_7e_7$ of the octonion $x$ (section 7.7) \\
$S(\beta),\ R(\theta),\ \Lambda$ & \texttt{Sβ}, \texttt{Rθ}, \texttt{Λboost}, \texttt{Λrot} & a boost, a rotation and their Lorentz matrices (section 8.2) \\
$\Gamma^{A}$ & \texttt{Γ17up[A]} & $\eta_{(1,7)}^{AA}\Gamma^{(1,7)}_A$, index raised with the inverse metric \\
$p,\ m,\ \not p$ & \texttt{pvec}, \texttt{mass}, \texttt{pslash} & $p=(3,2,2,0,0,0,0,0)$, $m=1$, $\not p=\sum_Ap^A\Gamma^{(1,7)}_A$ (section 8.4) \\
$W$ & \texttt{realAmplitudes} & the 8 real solutions $w$ of $(\not p-m)w=0$, one per row \\
$k,\ \not k$ & \texttt{kvec}, \texttt{kslash} & $k=(1,1,0,0,0,0,0,0)$, $\not k=\sum_Ak^A\Gamma^{(1,7)}_A$ (section 8.5) \\
$\widehat\Gamma^{A}\partial_A,\ P$ & \texttt{DiracOp71R32}, \texttt{pslash32} & $\sum_A\eta_{(7,1)}^{AA}\widehat\Gamma^{(7,1)}_A\partial_A$ and $P=\sum_Ap^A\widehat\Gamma^{(7,1)}_A$ (section 8.5) \\
$\Psi=(\Psi_1,\dots,\Psi_{16})^{\mathsf T}$ & \texttt{dirac16complex} & section 11: the quantized 16-component field, complex anticommuting (Grassmann-odd) operators; also, in section 8.3, 16 arbitrary functions \\
$\bar\Psi$, $\Psi^c$ & — & $\Psi^\dagger\Gamma^{0}$ (Dirac adjoint); $\Gamma^{(1,7)}_8\Psi^*$ (charge conjugate) \\
$\Psi_\uparrow$, $\Psi_\downarrow$ & — & upper and lower 8 components of $\Psi$ \\
$\Psi_L$, $\Psi_R$ & \texttt{ϕL}, \texttt{ϕR} & $(\Psi_\uparrow+i\Psi_\downarrow)/\sqrt2$, $(i\Psi_\uparrow+\Psi_\downarrow)/\sqrt2$, chiral 8-component fields \\
$P^{(1,7)}_L$, $P^{(1,7)}_R$ & \texttt{PL17}, \texttt{PR17} & $\frac12(1_{16}\mp\chi^{(1,7)})$, the complex chiral projections of 7+1 spacetime \\
$\chi$ & \texttt{χ17} & in section 11: $\chi=\chi^{(1,7)}=i\Gamma^{(1,7)}_8$ \\
$m$ & \texttt{mD}, \texttt{mass} & mass ($m=1$ in the examples) \\
$\vec p$, $E$ & \texttt{pv7}, \texttt{Ep} & momentum $(p^1,\dots,p^7)$; $E=\sqrt{m^2+\vec p^{\,2}}$ \\
$H_{\rm D}$ & \texttt{hD[ψ, m]} & $-i\Gamma^{0}\Gamma^{i}\partial_i+m\Gamma^{0}$, the Dirac Hamiltonian: $i\partial_0\Psi=H_{\rm D}\Psi$ \\
$e_s$, $w_s$ & \texttt{UnitVector[8, s]}, \texttt{uS[s, M, E]} & unit column of length 8; $w_s=\frac1{\sqrt E}(E\,e_s;\ (m1_8-p^i\tau^{O}_i)e_s)$, real positive-energy amplitude, $s=1,\dots,8$ \\
$g_S$, $g_P$, $g_V$, $g_A$ & \texttt{gS}, \texttt{gP}, \texttt{gV}, \texttt{gA} & real couplings of the scalar, pseudoscalar, vector and axial-vector contact interactions (mass dimension $-6$) \\
$\mathsf s$, $\mathsf p$, $\mathsf v_A$, $\mathsf a_A$ & — & real external sources \\
$\rho_S$, $\rho_P$, $j^A$, $j_5^A$ & \texttt{dens16[M]} & $\bar\Psi\Psi$, $\bar\Psi i\chi\Psi$, $\bar\Psi\Gamma^{A}\Psi$, $\bar\Psi\Gamma^{A}\chi\Psi$, or the sums over Kohn–Sham orbitals; $j_A=\eta^{(1,7)}_{AA}j^A$ \\
$\Phi_k$, $\lambda_k$ & \texttt{orbV} & Kohn–Sham orbitals (16 components) and their eigenvalues \\
$\Phi_{kL}$, $\Phi_{kR}$ & \texttt{orbLV}, \texttt{orbRV} & chiral components of $\Phi_k$, formed like $\Psi_L$, $\Psi_R$ \\
$H_{\rm KS}$ & \texttt{hKS[p]} & Kohn–Sham Hamiltonian (section 11.3) \\
$\mathcal S$, $\mathcal P$, $\mathcal V_A$, $\mathcal A_A$ & \texttt{Sks}, \texttt{Pks}, \texttt{VKS}, \texttt{AKS} & Kohn–Sham potentials (lower index $A$) \\
$F$, $T_s$, $E_H$, $E_{xc}$, $E_x$ & \texttt{EH16}, \texttt{Exfun} & universal functional; kinetic energy of the orbitals; Hartree, exchange-correlation and exchange energies \\
$\mathcal R$ & \texttt{Rorb} & $\sum_k\Phi_k\bar\Phi_k$, 16×16 \\
$\kappa$ & \texttt{κc} & $\sum_k\Phi_{kL}^\dagger\Phi_{kR}$, complex scalar density \\
$\mathcal M$, $\mathcal V^{L}_A$, $\mathcal V^{R}_A$ & \texttt{Mks}, \texttt{VLks}, \texttt{VRks} & $m+\mathcal S+i\mathcal P$ (complex Kohn–Sham mass); $\mathcal V_A-\mathcal A_A$; $\mathcal V_A+\mathcal A_A$ \\
$\vartheta$ & \texttt{Uθ} & angle of a chiral rotation $e^{i\vartheta\chi}$ \\
$k_F$, $m^*$, $\rho$, $E/A$, $E_F$ & \texttt{kF0}, \texttt{gap7}, \texttt{rho7}, \texttt{eps7}, \texttt{fermi7} & uniform matter: Fermi momentum, effective mass, density $j^0$, energy per particle, Fermi energy \\
$\Lambda_{\rm c}$, $g$, $g_c$ & \texttt{Λc}, \texttt{gNJL}, \texttt{gcrit} & momentum cutoff of the Dirac sea; $g=g_S=g_P$; critical coupling $-90\pi^4/\Lambda_{\rm c}^6$ \\
$g_{\mu\nu}$, $\{{}^{\lambda}_{\mu\nu}\}$ & \texttt{cv["g"]}, \texttt{cv["chr"]} & section 12: metric (one time-like and seven space-like dimensions, xAct signature $\{7,1,0\}$; components with the signs of $\eta^{(1,7)}$, $g_{00}>0$) and Christoffel symbols (never written $\Gamma$) \\
$R^\rho{}_{\sigma\mu\nu}$, $R_{\sigma\nu}$, $R$, $G^\mu{}_\nu$ & \texttt{cv["riem"]}, \texttt{cv["ric"]}, \texttt{cv["R"]}, \texttt{cv["G"]} & Riemann (Landau–Lifshitz), Ricci, scalar curvature, Einstein tensor \\
$E^{(k)\mu}{}_\nu$, $L_k$ & \texttt{Lovelock1/2/3}, \texttt{lovelockE}, \texttt{lovelockL} & Lovelock tensors and densities ($k=1,2,3$) \\
$\alpha_k$, $\bar\Lambda$, $\kappa$ & \texttt{alf[k]}, \texttt{lamb}, \texttt{kap} & Lovelock couplings ($\alpha_1=1$), cosmological constant, $8\pi G_8$ \\
$w_k$, $\Lambda$, $H$ & \texttt{w1}, \texttt{w2}, \texttt{w3}, \texttt{Λ}, \texttt{H} & pure-number couplings and inverse length of the reference notebook \\
$a$, $b$, $H_a$, $H_b$ & \texttt{a3[x0]}, \texttt{b4[x0]}, \texttt{Ha}, \texttt{Hb} & scale factors of ordinary and auxiliary space; Hubble rates $a'/a$, $b'/b$ \\
$e^A{}_\mu$, $\Gamma_\mu(X)$, $\Omega_\mu$ & \texttt{fd["e"]}, \texttt{fd["Gam"]}, \texttt{fd["Om"]} & frame field, generalized frame matrices, spinor connection \\
$\Theta^K$, $I_r$, $g_r$ & \texttt{Θup[K]}, \texttt{Ichan}, \texttt{gg[r]} & Clifford products of rank $r$, their invariants and couplings \\
$F_{rs}$, $G_s$ & \texttt{fierzF}, \texttt{Geff} & Fierz matrix and Hartree–Fock couplings $G=(1-F^{\mathsf T})g$ \\
$M$, $\varepsilon$, $p_a$, $p_b$ & \texttt{ksStress} & effective mass, energy density and pressures of Kohn–Sham matter \\
$W$ (brane), $h_{\mu\nu}$ & — & world brane $x_4=\dots=x_7=0$, induced metric \\
$v$, $V_{\rm CKM}$, $U_{\rm PMNS}$ & \texttt{vEW}, \texttt{vCKM}, \texttt{uPMNS} & Higgs vacuum value, quark and lepton mixing matrices \\
$R_A$, $U$ & \texttt{rBr}, \texttt{uInt} & brane form of the Dirac matrices and the intertwiner $\Gamma^{(1,7)}_AU=UR_A$ \\
$P(\alpha\to\beta)$ & \texttt{pVac}, \texttt{pMat} & neutrino flavour-conversion probability \\
$\beta$, $m_r$, $U(b)$ & \texttt{betR}, \texttt{rIIm}, \texttt{uPot} & radion $\ln(b/b_0)$, its mass $m_r^2=\kappa K_U/12$, the stabilizing potential (sections 20–23) \\
$u$, $X$, $P_E,P_G,P_C$ & \texttt{uS}, \texttt{xS}, \texttt{rIpE}, \texttt{rIpG}, \texttt{rIpC} & variables and polynomials of the exact primordial vacuum curve \\
$p_1,q_1$; $p_2,q_2$ & \texttt{rIp1}, \texttt{rIq1}; \texttt{rIp2}, \texttt{rIq2} & Einstein- and Gauss–Bonnet–Kasner exponents \\
$S_1,S_2,S_3$, $f_{\rm gas}$ & \texttt{rIstate[h]["S"]}, \texttt{["frac"]} & sizes of the curvature terms of the 00 equation; gas fraction \\
$g_*$, $g_{*s}$, $t(T)$ & \texttt{rIIgstars}, \texttt{rIItg} & degrees of freedom of the brane plasma, time–temperature relation \\
$H_\infty$, $\gamma$, $m_{\rm eff}$ & \texttt{rIIIH}, \texttt{rIIIgam}, \texttt{rIIImeff} & final de Sitter rate; radion friction and mass in de Sitter \\*
\end{longtable}}
\Needspace{12\baselineskip}\section*{4. The two definitions, and four corrections}\addcontentsline{toc}{section}{4. The two definitions, and four corrections}
\Needspace{8\baselineskip}\subsection*{4.1 [1A] The octonions}\addcontentsline{toc}{subsection}{4.1 [1A] The octonions}
The octonions $\mathbb{O}$ form an 8-dimensional, non-associative, non-commutative normed division algebra over $\mathbb{R}$, with basis $e_0=1,e_1,\dots,e_7$:
\begin{align*}
x&=\textstyle\sum_{A=0}^{7}x_A e_A,\qquad e_i^2=-1\ (i=1,\dots,7),\qquad e_ie_j=-e_je_i\ (i\neq j), \\
\bar x&=x_0-\textstyle\sum_{i=1}^{7}x_ie_i,\qquad \overline{xy}=\bar y\,\bar x,\qquad \mathrm{Re}(x)=\tfrac12(x+\bar x)=x_0, \\
N(x)&=x\bar x=\bar x x=x_0^2+x_1^2+\dots+x_7^2,\qquad N(xy)=N(x)N(y),\qquad x^{-1}=\bar x/N(x), \\
x(xy)&=(xx)y,\qquad (yx)x=y(xx)\qquad\text{(alternativity; power-associativity follows).}
\end{align*}
Products of distinct imaginary units follow the oriented lines of the Fano plane, each line spanning a quaternion subalgebra. Cayley–Dickson: $\mathbb{O}=\mathbb{H}\oplus\mathbb{H}$ as pairs of quaternions.

\Needspace{8\baselineskip}\subsection*{4.2 [1B] The Clifford algebra of an 8-dimensional space}\addcontentsline{toc}{subsection}{4.2 [1B] The Clifford algebra of an 8-dimensional space}
\begin{align*}
Cl(V,Q)&=\mathcal T(V)/\mathcal I_Q(V),\qquad \mathcal I_Q(V)=\big\langle v\otimes v-Q(v)1\big\rangle, \\
v^2&=Q(v)\,1,\qquad vw+wv=2B(v,w)\,1,\qquad B(v,w)=\tfrac12\big(Q(v+w)-Q(v)-Q(w)\big), \\
e_A^2&=Q(e_A),\qquad e_Ae_B=-e_Be_A\ (A\neq B),\qquad \dim Cl(V,Q)=2^8=\textstyle\sum_{k=0}^{8}\binom8k=256, \\
Cl_8(\mathbb C)&\cong\mathbb C(16).
\end{align*}
\Needspace{8\baselineskip}\subsection*{4.3 Corrections (each verified in 7+1.nb)}\addcontentsline{toc}{subsection}{4.3 Corrections (each verified in 7+1.nb)}
\begin{enumerate}
\item \textbf{Cayley–Dickson formula.} The second formula of [1A], $(a,b)(c,d)=(ac-\bar d b,\ \bar a d+cb)$, does not define the octonions: it is not alternative and its norm is not multiplicative. Counterexample: $x=(1,1)=e_0+e_4$, $y=(i,0)=e_1$ give $x^2=(0,2)$, $x^2y=(0,2i)$, but $xy=(i,i)$ and $x(xy)=(2i,2i)$. A correct formula of the same shape is $(a,b)(c,d)=(ac-d\bar b,\ \bar a d+cb)$ (J. C. Baez, \emph{The Octonions}, Bull. Amer. Math. Soc. 39 (2002), section 2.2); the octonions of this document satisfy it exactly with $x=a+\ell b$, $\ell=e_4$. The first formula of [1A] contains a stray symbol "x" and cannot be read unambiguously.
\item \textbf{Grade 4.} $\binom84=70$, not 72: $1+8+28+56+70+56+28+8+1=256$.
\item \textbf{Classification.} With $v^2=Q(v)$ and $p$ = number of generators squaring to $+1$, the correct table is the one below. $\mathbb C(16)$ can never occur: its real dimension is $2\cdot16^2=512$, while every $Cl(p,q)$ with $p+q=8$ has real dimension $2^8=256$. Several "Metric Signs" entries of [1B] were also mistyped.
\item \textbf{Split octonions are not a division algebra.} In the split-octonion table of 4+4.nb, $(\varepsilon_1+\varepsilon_5)(\varepsilon_1-\varepsilon_5)=\varepsilon_1-\varepsilon_5^2=0$ although both factors are nonzero. They form a composition algebra with an indefinite norm of signature (4,4).
\end{enumerate}
{\small
\begin{longtable}{>{\raggedright\arraybackslash}p{\dimexpr 0.1277\linewidth-2\tabcolsep\relax}>{\raggedright\arraybackslash}p{\dimexpr 0.2766\linewidth-2\tabcolsep\relax}>{\raggedright\arraybackslash}p{\dimexpr 0.1489\linewidth-2\tabcolsep\relax}>{\raggedright\arraybackslash}p{\dimexpr 0.1915\linewidth-2\tabcolsep\relax}>{\raggedright\arraybackslash}p{\dimexpr 0.2553\linewidth-2\tabcolsep\relax}}
\toprule
\textbf{$Cl(p,q)$} & \textbf{squares of the generators} & \textbf{correct} & \textbf{claimed in [1B]} & \textbf{status} \\
\midrule\endfirsthead
\toprule
\textbf{$Cl(p,q)$} & \textbf{squares of the generators} & \textbf{correct} & \textbf{claimed in [1B]} & \textbf{status} \\
\midrule\endhead
\endfoot
\bottomrule\endlastfoot
$Cl(8,0)$ & +,+,+,+,+,+,+,+ & $\mathbb R(16)$ & $\mathbb R(16)$ & agrees \\
$Cl(7,1)$ & +,+,+,+,+,+,+,− & $\mathbb H(8)$ & $\mathbb R(16)$ & wrong in [1B] \\
$Cl(6,2)$ & +,+,+,+,+,+,−,− & $\mathbb H(8)$ & $\mathbb H(8)$ & agrees (signs mistyped) \\
$Cl(5,3)$ & +,+,+,+,+,−,−,− & $\mathbb R(16)$ & $\mathbb C(16)$ & wrong in [1B] \\
$Cl(4,4)$ & +,+,+,+,−,−,−,− & $\mathbb R(16)$ & $\mathbb R(16)$ & agrees (signs mistyped) \\
$Cl(3,5)$ & +,+,+,−,−,−,−,− & $\mathbb H(8)$ & $\mathbb C(16)$ & wrong in [1B] \\
$Cl(2,6)$ & +,+,−,−,−,−,−,− & $\mathbb H(8)$ & $\mathbb H(8)$ & agrees (signs mistyped) \\
$Cl(1,7)$ & +,−,−,−,−,−,−,− & $\mathbb R(16)$ & $\mathbb R(16)$ & agrees \\
$Cl(0,8)$ & −,−,−,−,−,−,−,− & $\mathbb R(16)$ & $\mathbb R(16)$ & agrees \\*
\end{longtable}}
\textbf{Consequence for this project.} With $v^2=Q(v)$ and the sign convention of 4+4.nb ($Q=\eta^{(7,1)}$) the Clifford algebra of 7+1 spacetime is $Cl(7,1)\cong\mathbb H(8)$: no real 16×16 matrices exist (section 7.4). With the opposite overall sign ($Q=\eta^{(1,7)}$, the octonions' own form $\mathrm{Re}(xx)$) it is $Cl(1,7)\cong\mathbb R(16)$, realized by the real matrices $\Gamma^{(1,7)}_A$ built from 4+4.nb (section 7.2). Both describe the same spacetime: one time and seven space dimensions.

\Needspace{12\baselineskip}\section*{5. Recap: the 4+4 construction, reproduced exactly from 4+4.nb}\addcontentsline{toc}{section}{5. Recap: the 4+4 construction, reproduced exactly from 4+4.nb}
\Needspace{8\baselineskip}\subsection*{5.1 The 8×8 matrices \texorpdfstring{$\sigma$}{σ}, \texorpdfstring{$\zeta_A$}{ζA}, \texorpdfstring{$\tau_A$}{τA}, \texorpdfstring{$\bar\tau_A$}{τ̄A}}\addcontentsline{toc}{subsection}{5.1 The 8×8 matrices \texorpdfstring{$\sigma$}{σ}, \texorpdfstring{$\zeta_A$}{ζA}, \texorpdfstring{$\tau_A$}{τA}, \texorpdfstring{$\bar\tau_A$}{τ̄A}}
With $\epsilon_{hpq4}$ the sign of the permutation $(h,p,q,4)$ (Mathematica \texttt{Signature[\{h, p, q, 4\}]}), 4+4.nb defines the self-dual and anti-self-dual antisymmetric 4×4 matrices
\begin{equation*}
(s_h)_{pq}=\epsilon_{hpq4}-\big(\delta_{p4}\delta_{qh}-\delta_{ph}\delta_{q4}\big),\qquad (t_h)_{pq}=\epsilon_{hpq4}+\big(\delta_{p4}\delta_{qh}-\delta_{ph}\delta_{q4}\big),\qquad h=1,2,3,\ p,q=1,\dots,4,
\end{equation*}
\begin{equation*}
s_1=\begin{pmatrix}0 & 0 & 0 & 1 \\ 0 & 0 & 1 & 0 \\ 0 & -1 & 0 & 0 \\ -1 & 0 & 0 & 0\end{pmatrix},\ s_2=\begin{pmatrix}0 & 0 & -1 & 0 \\ 0 & 0 & 0 & 1 \\ 1 & 0 & 0 & 0 \\ 0 & -1 & 0 & 0\end{pmatrix},\ s_3=\begin{pmatrix}0 & 1 & 0 & 0 \\ -1 & 0 & 0 & 0 \\ 0 & 0 & 0 & 1 \\ 0 & 0 & -1 & 0\end{pmatrix},
\end{equation*}
\begin{equation*}
t_1=\begin{pmatrix}0 & 0 & 0 & -1 \\ 0 & 0 & 1 & 0 \\ 0 & -1 & 0 & 0 \\ 1 & 0 & 0 & 0\end{pmatrix},\ t_2=\begin{pmatrix}0 & 0 & -1 & 0 \\ 0 & 0 & 0 & -1 \\ 1 & 0 & 0 & 0 \\ 0 & 1 & 0 & 0\end{pmatrix},\ t_3=\begin{pmatrix}0 & 1 & 0 & 0 \\ -1 & 0 & 0 & 0 \\ 0 & 0 & 0 & -1 \\ 0 & 0 & 1 & 0\end{pmatrix},
\end{equation*}
and from them the 8×8 matrices
\begin{align*}
\sigma&=\begin{pmatrix}0&1_4\\1_4&0\end{pmatrix},\qquad \zeta_0=1_8,\qquad \zeta_h=\begin{pmatrix}0&s_h\\s_h&0\end{pmatrix},\qquad \zeta_{7-h}=\begin{pmatrix}0&t_h\\-t_h&0\end{pmatrix}\ (h=1,2,3), \\
\zeta_7&=\zeta_1\zeta_2\zeta_3\zeta_4\zeta_5\zeta_6=\mathrm{diag}(-1_4,1_4),\qquad \tau_A=\zeta_A,\qquad \bar\tau_0=1_8,\quad \bar\tau_A=\sigma\,\tau_A^{\mathsf T}\sigma\ (A=1,\dots,7).
\end{align*}
\begin{equation*}
\tau_1=\begin{pmatrix}0 & 0 & 0 & 0 & 0 & 0 & 0 & 1 \\ 0 & 0 & 0 & 0 & 0 & 0 & 1 & 0 \\ 0 & 0 & 0 & 0 & 0 & -1 & 0 & 0 \\ 0 & 0 & 0 & 0 & -1 & 0 & 0 & 0 \\ 0 & 0 & 0 & 1 & 0 & 0 & 0 & 0 \\ 0 & 0 & 1 & 0 & 0 & 0 & 0 & 0 \\ 0 & -1 & 0 & 0 & 0 & 0 & 0 & 0 \\ -1 & 0 & 0 & 0 & 0 & 0 & 0 & 0\end{pmatrix},\quad \tau_2=\begin{pmatrix}0 & 0 & 0 & 0 & 0 & 0 & -1 & 0 \\ 0 & 0 & 0 & 0 & 0 & 0 & 0 & 1 \\ 0 & 0 & 0 & 0 & 1 & 0 & 0 & 0 \\ 0 & 0 & 0 & 0 & 0 & -1 & 0 & 0 \\ 0 & 0 & -1 & 0 & 0 & 0 & 0 & 0 \\ 0 & 0 & 0 & 1 & 0 & 0 & 0 & 0 \\ 1 & 0 & 0 & 0 & 0 & 0 & 0 & 0 \\ 0 & -1 & 0 & 0 & 0 & 0 & 0 & 0\end{pmatrix},
\end{equation*}
\begin{equation*}
\tau_3=\begin{pmatrix}0 & 0 & 0 & 0 & 0 & 1 & 0 & 0 \\ 0 & 0 & 0 & 0 & -1 & 0 & 0 & 0 \\ 0 & 0 & 0 & 0 & 0 & 0 & 0 & 1 \\ 0 & 0 & 0 & 0 & 0 & 0 & -1 & 0 \\ 0 & 1 & 0 & 0 & 0 & 0 & 0 & 0 \\ -1 & 0 & 0 & 0 & 0 & 0 & 0 & 0 \\ 0 & 0 & 0 & 1 & 0 & 0 & 0 & 0 \\ 0 & 0 & -1 & 0 & 0 & 0 & 0 & 0\end{pmatrix},\quad \tau_4=\begin{pmatrix}0 & 0 & 0 & 0 & 0 & 1 & 0 & 0 \\ 0 & 0 & 0 & 0 & -1 & 0 & 0 & 0 \\ 0 & 0 & 0 & 0 & 0 & 0 & 0 & -1 \\ 0 & 0 & 0 & 0 & 0 & 0 & 1 & 0 \\ 0 & -1 & 0 & 0 & 0 & 0 & 0 & 0 \\ 1 & 0 & 0 & 0 & 0 & 0 & 0 & 0 \\ 0 & 0 & 0 & 1 & 0 & 0 & 0 & 0 \\ 0 & 0 & -1 & 0 & 0 & 0 & 0 & 0\end{pmatrix},
\end{equation*}
\begin{equation*}
\tau_5=\begin{pmatrix}0 & 0 & 0 & 0 & 0 & 0 & -1 & 0 \\ 0 & 0 & 0 & 0 & 0 & 0 & 0 & -1 \\ 0 & 0 & 0 & 0 & 1 & 0 & 0 & 0 \\ 0 & 0 & 0 & 0 & 0 & 1 & 0 & 0 \\ 0 & 0 & 1 & 0 & 0 & 0 & 0 & 0 \\ 0 & 0 & 0 & 1 & 0 & 0 & 0 & 0 \\ -1 & 0 & 0 & 0 & 0 & 0 & 0 & 0 \\ 0 & -1 & 0 & 0 & 0 & 0 & 0 & 0\end{pmatrix},\quad \tau_6=\begin{pmatrix}0 & 0 & 0 & 0 & 0 & 0 & 0 & -1 \\ 0 & 0 & 0 & 0 & 0 & 0 & 1 & 0 \\ 0 & 0 & 0 & 0 & 0 & -1 & 0 & 0 \\ 0 & 0 & 0 & 0 & 1 & 0 & 0 & 0 \\ 0 & 0 & 0 & 1 & 0 & 0 & 0 & 0 \\ 0 & 0 & -1 & 0 & 0 & 0 & 0 & 0 \\ 0 & 1 & 0 & 0 & 0 & 0 & 0 & 0 \\ -1 & 0 & 0 & 0 & 0 & 0 & 0 & 0\end{pmatrix},
\end{equation*}
\begin{equation*}
\tau_7=\begin{pmatrix}-1 & 0 & 0 & 0 & 0 & 0 & 0 & 0 \\ 0 & -1 & 0 & 0 & 0 & 0 & 0 & 0 \\ 0 & 0 & -1 & 0 & 0 & 0 & 0 & 0 \\ 0 & 0 & 0 & -1 & 0 & 0 & 0 & 0 \\ 0 & 0 & 0 & 0 & 1 & 0 & 0 & 0 \\ 0 & 0 & 0 & 0 & 0 & 1 & 0 & 0 \\ 0 & 0 & 0 & 0 & 0 & 0 & 1 & 0 \\ 0 & 0 & 0 & 0 & 0 & 0 & 0 & 1\end{pmatrix}.
\end{equation*}
They satisfy (all verified in 7+1.nb exactly as in 4+4.nb)
\begin{align*}
\tfrac12\big(\zeta_A\zeta_B+\zeta_B\zeta_A\big)&=-\eta^{(4,4)}_{AB}\,1_8\quad(A,B=1,\dots,7),\qquad -\eta^{(4,4)}_{AA}\,\zeta_A=\zeta_A^{\mathsf T}, \\
\tfrac12\big(\tau_A\bar\tau_B+\tau_B\bar\tau_A\big)&=\tfrac12\big(\bar\tau_A\tau_B+\bar\tau_B\tau_A\big)=\eta^{(4,4)}_{AB}\,1_8\quad(A,B=0,\dots,7), \\
\tau_1\tau_2\cdots\tau_7&=1_8,\qquad \sigma=\tau_1\tau_2\tau_3=\tau_4\tau_5\tau_6\tau_7,\qquad \sigma^2=1_8, \\
\sigma\tau_h&=\tau_h\sigma\ (h=1,2,3),\qquad \sigma\tau_k=-\tau_k\sigma\ (k=4,\dots,7).
\end{align*}
So $\tau_1,\tau_2,\tau_3$ are antisymmetric and square to $-1_8$; $\tau_4,\dots,\tau_7$ are symmetric and square to $+1_8$: they are the left multiplications of the three imaginary and the four split units of the split octonions.

\Needspace{8\baselineskip}\subsection*{5.2 The 16×16 Dirac matrices \texorpdfstring{$T_A$}{TA} of 4+4.nb}\addcontentsline{toc}{subsection}{5.2 The 16×16 Dirac matrices \texorpdfstring{$T_A$}{TA} of 4+4.nb}
\begin{align*}
T_A&=\begin{pmatrix}0&\bar\tau_A\\ \tau_A&0\end{pmatrix}\ (A=0,\dots,7),\qquad T_AT_B+T_BT_A=2\,\eta^{(4,4)}_{AB}\,1_{16}, \\
T_8&=T_0T_1\cdots T_7=\mathrm{diag}(-1_8,1_8),\qquad \sigma_{16}=T_0T_1T_2T_3=\mathrm{diag}(-\sigma,\sigma), \\
P_L&=\tfrac12(1_{16}-T_8)=\mathrm{diag}(1_8,0),\qquad P_R=\tfrac12(1_{16}+T_8)=\mathrm{diag}(0,1_8).
\end{align*}
\Needspace{8\baselineskip}\subsection*{5.3 How the split-octonion multiplication constants come from \texorpdfstring{$\tau_A$}{τA}}\addcontentsline{toc}{subsection}{5.3 How the split-octonion multiplication constants come from \texorpdfstring{$\tau_A$}{τA}}
Choose the unit spinor $u=(1,0,0,0,1,0,0,0)^{\mathsf T}/\sqrt2$, an eigenvector of $\sigma$: $\sigma u=u$, $u^{\mathsf T}u=u^{\mathsf T}\sigma u=1$. The vielbein $E_{44}$ (4+4.nb: \texttt{E\textasciicircum{}A\_a[hUSE]}) and its inverse are
\begin{equation*}
(E_{44})_{Ba}=\big(u^{\mathsf T}\sigma\,\bar\tau_B\big)_a,\qquad E_{44}^{-1}=(u^{\mathsf T}\sigma u)\,\big(\eta^{(4,4)}_{00}\tau_0u,\ \eta^{(4,4)}_{11}\tau_1u,\ \dots,\ \eta^{(4,4)}_{77}\tau_7u\big)\quad\text{(columns)}.
\end{equation*}
The left multiplication by the split unit $\varepsilon_{B}$ (4+4.nb: \texttt{eA[B]}, $B=1,\dots,8$) is the matrix $g_{BB}\,E_{44}\,\tau_{B-1}\,E_{44}^{-1}$ with $g=\eta^{(4,4)}$:
\begin{equation*}
\varepsilon_B\,\varepsilon_C=\sum_{D=1}^{8}\Big(g_{BB}\,E_{44}\,\tau_{B-1}\,E_{44}^{-1}\Big)_{DC}\,\varepsilon_D .
\end{equation*}
In words: the split-octonion multiplication constants are the matrix elements of the 8×8 matrices $\tau_A$, written in the octonion basis by the vielbein, times the metric sign $g_{BB}$. This is the table \texttt{gridx} of 4+4.nb (row $\varepsilon_B$ times column $\varepsilon_C$):

{\footnotesize
\begin{longtable}{>{\raggedright\arraybackslash}p{\dimexpr 0.1064\linewidth-2\tabcolsep\relax}>{\raggedright\arraybackslash}p{\dimexpr 0.1117\linewidth-2\tabcolsep\relax}>{\raggedright\arraybackslash}p{\dimexpr 0.1117\linewidth-2\tabcolsep\relax}>{\raggedright\arraybackslash}p{\dimexpr 0.1117\linewidth-2\tabcolsep\relax}>{\raggedright\arraybackslash}p{\dimexpr 0.1117\linewidth-2\tabcolsep\relax}>{\raggedright\arraybackslash}p{\dimexpr 0.1117\linewidth-2\tabcolsep\relax}>{\raggedright\arraybackslash}p{\dimexpr 0.1117\linewidth-2\tabcolsep\relax}>{\raggedright\arraybackslash}p{\dimexpr 0.1117\linewidth-2\tabcolsep\relax}>{\raggedright\arraybackslash}p{\dimexpr 0.1117\linewidth-2\tabcolsep\relax}}
\toprule
\textbf{$\cdot$} & \textbf{$\varepsilon_{1}$} & \textbf{$\varepsilon_{2}$} & \textbf{$\varepsilon_{3}$} & \textbf{$\varepsilon_{4}$} & \textbf{$\varepsilon_{5}$} & \textbf{$\varepsilon_{6}$} & \textbf{$\varepsilon_{7}$} & \textbf{$\varepsilon_{8}$} \\
\midrule\endfirsthead
\toprule
\textbf{$\cdot$} & \textbf{$\varepsilon_{1}$} & \textbf{$\varepsilon_{2}$} & \textbf{$\varepsilon_{3}$} & \textbf{$\varepsilon_{4}$} & \textbf{$\varepsilon_{5}$} & \textbf{$\varepsilon_{6}$} & \textbf{$\varepsilon_{7}$} & \textbf{$\varepsilon_{8}$} \\
\midrule\endhead
\endfoot
\bottomrule\endlastfoot
$\varepsilon_{1}$ & $\varepsilon_{1}$ & $\varepsilon_{2}$ & $\varepsilon_{3}$ & $\varepsilon_{4}$ & $\varepsilon_{5}$ & $\varepsilon_{6}$ & $\varepsilon_{7}$ & $\varepsilon_{8}$ \\
$\varepsilon_{2}$ & $\varepsilon_{2}$ & $-\varepsilon_{1}$ & $-\varepsilon_{4}$ & $\varepsilon_{3}$ & $-\varepsilon_{6}$ & $\varepsilon_{5}$ & $\varepsilon_{8}$ & $-\varepsilon_{7}$ \\
$\varepsilon_{3}$ & $\varepsilon_{3}$ & $\varepsilon_{4}$ & $-\varepsilon_{1}$ & $-\varepsilon_{2}$ & $-\varepsilon_{7}$ & $-\varepsilon_{8}$ & $\varepsilon_{5}$ & $\varepsilon_{6}$ \\
$\varepsilon_{4}$ & $\varepsilon_{4}$ & $-\varepsilon_{3}$ & $\varepsilon_{2}$ & $-\varepsilon_{1}$ & $-\varepsilon_{8}$ & $\varepsilon_{7}$ & $-\varepsilon_{6}$ & $\varepsilon_{5}$ \\
$\varepsilon_{5}$ & $\varepsilon_{5}$ & $\varepsilon_{6}$ & $\varepsilon_{7}$ & $\varepsilon_{8}$ & $\varepsilon_{1}$ & $\varepsilon_{2}$ & $\varepsilon_{3}$ & $\varepsilon_{4}$ \\
$\varepsilon_{6}$ & $\varepsilon_{6}$ & $-\varepsilon_{5}$ & $\varepsilon_{8}$ & $-\varepsilon_{7}$ & $-\varepsilon_{2}$ & $\varepsilon_{1}$ & $-\varepsilon_{4}$ & $\varepsilon_{3}$ \\
$\varepsilon_{7}$ & $\varepsilon_{7}$ & $-\varepsilon_{8}$ & $-\varepsilon_{5}$ & $\varepsilon_{6}$ & $-\varepsilon_{3}$ & $\varepsilon_{4}$ & $\varepsilon_{1}$ & $-\varepsilon_{2}$ \\
$\varepsilon_{8}$ & $\varepsilon_{8}$ & $\varepsilon_{7}$ & $-\varepsilon_{6}$ & $-\varepsilon_{5}$ & $-\varepsilon_{4}$ & $-\varepsilon_{3}$ & $\varepsilon_{2}$ & $\varepsilon_{1}$ \\*
\end{longtable}}
Split octonions: $\varepsilon_1=1$, $\varepsilon_2^2=\varepsilon_3^2=\varepsilon_4^2=-1$, $\varepsilon_5^2=\dots=\varepsilon_8^2=+1$, norm $N_S(x)=x^{\mathsf T}\eta^{(4,4)}x$ with $N_S(xy)=N_S(x)N_S(y)$, alternative, and with zero divisors: $(\varepsilon_1+\varepsilon_5)(\varepsilon_1-\varepsilon_5)=0$.

\Needspace{12\baselineskip}\section*{6. The octonions: a real realization mirroring the split octonions}\addcontentsline{toc}{section}{6. The octonions: a real realization mirroring the split octonions}
\Needspace{8\baselineskip}\subsection*{6.1 From \texorpdfstring{$\tau$}{τ} to \texorpdfstring{$\tau^{O}$}{τO}}\addcontentsline{toc}{subsection}{6.1 From \texorpdfstring{$\tau$}{τ} to \texorpdfstring{$\tau^{O}$}{τO}}
All seven imaginary octonion units must square to $-1$, but $\tau_4,\dots,\tau_7$ square to $+1$. Because $\sigma$ commutes with $\tau_1,\tau_2,\tau_3$, anticommutes with $\tau_4,\dots,\tau_7$ and squares to $1_8$,
\begin{equation*}
(\tau_k\sigma)(\tau_k\sigma)=-\tau_k\tau_k\,\sigma\sigma=-1_8\qquad(k=4,\dots,7),
\end{equation*}
and all anticommutation relations survive. Over $\mathbb C$ one would multiply $\tau_k$ by $i$; the real matrix $\sigma$ replaces that $i$. Definition:
\begin{equation*}
\tau^{O}_0=1_8,\qquad \tau^{O}_A=\tau_A\ (A=1,2,3),\qquad \tau^{O}_A=\tau_A\,\sigma\ (A=4,\dots,7),\qquad \bar\tau^{O}_A=g^{O}(\tau^{O}_A)^{\mathsf T}g^{O}=(\tau^{O}_A)^{\mathsf T}.
\end{equation*}
\begin{equation*}
\tau^{O}_4=\begin{pmatrix}0 & 1 & 0 & 0 & 0 & 0 & 0 & 0 \\ -1 & 0 & 0 & 0 & 0 & 0 & 0 & 0 \\ 0 & 0 & 0 & -1 & 0 & 0 & 0 & 0 \\ 0 & 0 & 1 & 0 & 0 & 0 & 0 & 0 \\ 0 & 0 & 0 & 0 & 0 & -1 & 0 & 0 \\ 0 & 0 & 0 & 0 & 1 & 0 & 0 & 0 \\ 0 & 0 & 0 & 0 & 0 & 0 & 0 & 1 \\ 0 & 0 & 0 & 0 & 0 & 0 & -1 & 0\end{pmatrix},\quad \tau^{O}_5=\begin{pmatrix}0 & 0 & -1 & 0 & 0 & 0 & 0 & 0 \\ 0 & 0 & 0 & -1 & 0 & 0 & 0 & 0 \\ 1 & 0 & 0 & 0 & 0 & 0 & 0 & 0 \\ 0 & 1 & 0 & 0 & 0 & 0 & 0 & 0 \\ 0 & 0 & 0 & 0 & 0 & 0 & 1 & 0 \\ 0 & 0 & 0 & 0 & 0 & 0 & 0 & 1 \\ 0 & 0 & 0 & 0 & -1 & 0 & 0 & 0 \\ 0 & 0 & 0 & 0 & 0 & -1 & 0 & 0\end{pmatrix},
\end{equation*}
\begin{equation*}
\tau^{O}_6=\begin{pmatrix}0 & 0 & 0 & -1 & 0 & 0 & 0 & 0 \\ 0 & 0 & 1 & 0 & 0 & 0 & 0 & 0 \\ 0 & -1 & 0 & 0 & 0 & 0 & 0 & 0 \\ 1 & 0 & 0 & 0 & 0 & 0 & 0 & 0 \\ 0 & 0 & 0 & 0 & 0 & 0 & 0 & 1 \\ 0 & 0 & 0 & 0 & 0 & 0 & -1 & 0 \\ 0 & 0 & 0 & 0 & 0 & 1 & 0 & 0 \\ 0 & 0 & 0 & 0 & -1 & 0 & 0 & 0\end{pmatrix},\quad \tau^{O}_7=\begin{pmatrix}0 & 0 & 0 & 0 & -1 & 0 & 0 & 0 \\ 0 & 0 & 0 & 0 & 0 & -1 & 0 & 0 \\ 0 & 0 & 0 & 0 & 0 & 0 & -1 & 0 \\ 0 & 0 & 0 & 0 & 0 & 0 & 0 & -1 \\ 1 & 0 & 0 & 0 & 0 & 0 & 0 & 0 \\ 0 & 1 & 0 & 0 & 0 & 0 & 0 & 0 \\ 0 & 0 & 1 & 0 & 0 & 0 & 0 & 0 \\ 0 & 0 & 0 & 1 & 0 & 0 & 0 & 0\end{pmatrix},
\end{equation*}
and $\tau^{O}_1=\tau_1$, $\tau^{O}_2=\tau_2$, $\tau^{O}_3=\tau_3$ as displayed in section 5.1. Properties (verified):
\begin{align*}
(\tau^{O}_i)^{\mathsf T}&=-\tau^{O}_i,\qquad (\tau^{O}_i)^2=-1_8,\qquad \tau^{O}_i\tau^{O}_j=-\tau^{O}_j\tau^{O}_i\ (i\neq j), \\
\tau^{O}_A\bar\tau^{O}_B+\tau^{O}_B\bar\tau^{O}_A&=\bar\tau^{O}_A\tau^{O}_B+\bar\tau^{O}_B\tau^{O}_A=2\,g^{O}_{AB}\,1_8=2\,\delta_{AB}\,1_8, \\
\tau^{O}_1\tau^{O}_2\cdots\tau^{O}_7&=1_8\qquad(\text{as }\tau_1\cdots\tau_7=1_8\text{ in 4+4.nb}), \\
\tau^{O}_4&=\tau_5\tau_6\tau_7,\quad \tau^{O}_5=-\tau_4\tau_6\tau_7,\quad \tau^{O}_6=\tau_4\tau_5\tau_7,\quad \tau^{O}_7=-\tau_4\tau_5\tau_6.
\end{align*}
For $x=\sum_Ax_Ae_A$ put $\tau^{O}(x)=\sum_Ax_A\tau^{O}_A$; then $\tau^{O}(x)\,\bar\tau^{O}(x)=N(x)\,1_8$ with $N(x)=x_0^2+\dots+x_7^2$, the matrix form of $x\bar x=N(x)$. In 4+4.nb the same relation holds with $\eta^{(4,4)}$, the split norm.

\Needspace{8\baselineskip}\subsection*{6.2 The octonion vielbein (same unit spinor \texorpdfstring{$u$}{u})}\addcontentsline{toc}{subsection}{6.2 The octonion vielbein (same unit spinor \texorpdfstring{$u$}{u})}
In 4+4.nb the rows of the vielbein are $u^{\mathsf T}\sigma\bar\tau_B$, with the split spinor metric $\sigma$. For the octonions the spinor metric is Euclidean ($\sigma\to g^{O}=1_8$), so
\begin{equation*}
(E^{O})_{Ba}=\big(u^{\mathsf T}g^{O}\bar\tau^{O}_B\big)_a,\qquad (E^{O})^{-1}=(E^{O})^{\mathsf T}=(u^{\mathsf T}g^{O}u)\big(g^{O}_{00}\tau^{O}_0u,\dots,g^{O}_{77}\tau^{O}_7u\big),\qquad E^{O}=\eta^{(4,4)}E_{44}.
\end{equation*}
Because $\sigma u=u$, $\tau^{O}_Au=\tau_Au$ for every $A$: the two vielbeins differ only by the signs of rows 5–8.
\begin{equation*}
E^{O}=\begin{pmatrix}\frac{\sqrt{2}}{2} & 0 & 0 & 0 & \frac{\sqrt{2}}{2} & 0 & 0 & 0 \\ 0 & 0 & 0 & - \frac{\sqrt{2}}{2} & 0 & 0 & 0 & - \frac{\sqrt{2}}{2} \\ 0 & 0 & \frac{\sqrt{2}}{2} & 0 & 0 & 0 & \frac{\sqrt{2}}{2} & 0 \\ 0 & - \frac{\sqrt{2}}{2} & 0 & 0 & 0 & - \frac{\sqrt{2}}{2} & 0 & 0 \\ 0 & - \frac{\sqrt{2}}{2} & 0 & 0 & 0 & \frac{\sqrt{2}}{2} & 0 & 0 \\ 0 & 0 & \frac{\sqrt{2}}{2} & 0 & 0 & 0 & - \frac{\sqrt{2}}{2} & 0 \\ 0 & 0 & 0 & \frac{\sqrt{2}}{2} & 0 & 0 & 0 & - \frac{\sqrt{2}}{2} \\ - \frac{\sqrt{2}}{2} & 0 & 0 & 0 & \frac{\sqrt{2}}{2} & 0 & 0 & 0\end{pmatrix} .
\end{equation*}
\Needspace{8\baselineskip}\subsection*{6.3 The octonion multiplication constants}\addcontentsline{toc}{subsection}{6.3 The octonion multiplication constants}
Exactly the formula of 4+4.nb, with $\tau\to\tau^{O}$, $E_{44}\to E^{O}$, $g\to g^{O}$ and a 0-based basis $e_A\leftrightarrow\tau^{O}_A\leftrightarrow x_A$:
\begin{equation*}
L^{O}_A=g^{O}_{AA}\,E^{O}\,\tau^{O}_A\,(E^{O})^{-1},\qquad e_A\,e_K=\sum_{L=0}^{7}\big(L^{O}_A\big)_{LK}\,e_L,\qquad (xy)_L=\sum_{A,K}x_A\,y_K\,\big(L^{O}_A\big)_{LK}.
\end{equation*}
Every $L^{O}_A$ is a signed permutation matrix, so every multiplication constant is $0$ or $\pm1$. The table (row $e_A$ times column $e_K$; Mathematica \texttt{gridO}, the mirror of \texttt{gridx}):

{\footnotesize
\begin{longtable}{>{\raggedright\arraybackslash}p{\dimexpr 0.1064\linewidth-2\tabcolsep\relax}>{\raggedright\arraybackslash}p{\dimexpr 0.1117\linewidth-2\tabcolsep\relax}>{\raggedright\arraybackslash}p{\dimexpr 0.1117\linewidth-2\tabcolsep\relax}>{\raggedright\arraybackslash}p{\dimexpr 0.1117\linewidth-2\tabcolsep\relax}>{\raggedright\arraybackslash}p{\dimexpr 0.1117\linewidth-2\tabcolsep\relax}>{\raggedright\arraybackslash}p{\dimexpr 0.1117\linewidth-2\tabcolsep\relax}>{\raggedright\arraybackslash}p{\dimexpr 0.1117\linewidth-2\tabcolsep\relax}>{\raggedright\arraybackslash}p{\dimexpr 0.1117\linewidth-2\tabcolsep\relax}>{\raggedright\arraybackslash}p{\dimexpr 0.1117\linewidth-2\tabcolsep\relax}}
\toprule
\textbf{$\cdot$} & \textbf{$e_{0}$} & \textbf{$e_{1}$} & \textbf{$e_{2}$} & \textbf{$e_{3}$} & \textbf{$e_{4}$} & \textbf{$e_{5}$} & \textbf{$e_{6}$} & \textbf{$e_{7}$} \\
\midrule\endfirsthead
\toprule
\textbf{$\cdot$} & \textbf{$e_{0}$} & \textbf{$e_{1}$} & \textbf{$e_{2}$} & \textbf{$e_{3}$} & \textbf{$e_{4}$} & \textbf{$e_{5}$} & \textbf{$e_{6}$} & \textbf{$e_{7}$} \\
\midrule\endhead
\endfoot
\bottomrule\endlastfoot
$e_{0}$ & $e_{0}$ & $e_{1}$ & $e_{2}$ & $e_{3}$ & $e_{4}$ & $e_{5}$ & $e_{6}$ & $e_{7}$ \\
$e_{1}$ & $e_{1}$ & $-e_{0}$ & $-e_{3}$ & $e_{2}$ & $-e_{5}$ & $e_{4}$ & $e_{7}$ & $-e_{6}$ \\
$e_{2}$ & $e_{2}$ & $e_{3}$ & $-e_{0}$ & $-e_{1}$ & $-e_{6}$ & $-e_{7}$ & $e_{4}$ & $e_{5}$ \\
$e_{3}$ & $e_{3}$ & $-e_{2}$ & $e_{1}$ & $-e_{0}$ & $-e_{7}$ & $e_{6}$ & $-e_{5}$ & $e_{4}$ \\
$e_{4}$ & $e_{4}$ & $e_{5}$ & $e_{6}$ & $e_{7}$ & $-e_{0}$ & $-e_{1}$ & $-e_{2}$ & $-e_{3}$ \\
$e_{5}$ & $e_{5}$ & $-e_{4}$ & $e_{7}$ & $-e_{6}$ & $e_{1}$ & $-e_{0}$ & $e_{3}$ & $-e_{2}$ \\
$e_{6}$ & $e_{6}$ & $-e_{7}$ & $-e_{4}$ & $e_{5}$ & $e_{2}$ & $-e_{3}$ & $-e_{0}$ & $e_{1}$ \\
$e_{7}$ & $e_{7}$ & $e_{6}$ & $-e_{5}$ & $-e_{4}$ & $e_{3}$ & $e_{2}$ & $-e_{1}$ & $-e_{0}$ \\*
\end{longtable}}
\Needspace{8\baselineskip}\subsection*{6.4 The properties of [1A], verified}\addcontentsline{toc}{subsection}{6.4 The properties of [1A], verified}
\textbf{Units.} $e_0$ is a two-sided identity, $e_i^2=-e_0$ and $e_ie_j=-e_je_i$ for $i\neq j$ ($i,j=1,\dots,7$).

\textbf{Fano plane.} Each product of two distinct imaginary units is $\pm$ a third one. The seven triples are the lines of the Fano plane; written in the cyclic order $(a,b,c)$ with $e_ae_b=+e_c$ (and then $e_be_c=e_a$, $e_ce_a=e_b$):
\begin{equation*}
(1,3,2),\ (1,5,4),\ (1,6,7),\ (2,6,4),\ (2,7,5),\ (3,5,6),\ (3,7,4).
\end{equation*}
Every pair of distinct imaginary units lies on exactly one line, and each line together with $e_0$ spans a copy of the quaternions (closed and associative). In a picture: put $e_1$ at the centre, $e_2,e_4,e_7$ at the vertices of a triangle and $e_6,e_3,e_5$ at the midpoints of the sides $e_2e_4$, $e_4e_7$, $e_7e_2$; the lines are the three sides, the three medians and the inscribed circle through $e_3,e_5,e_6$.

\textbf{Algebraic laws} (proved symbolically for general octonions $x,y,z$ with 8 symbolic components each; the Moufang identities, being linear in $x$ and $y$, are proved for basis units $x,y$ and symbolic $z$):
\begin{align*}
x(xy)&=(xx)y,\qquad (yx)x=y(xx),\qquad (xy)x=x(yx),\qquad x(xx)=(xx)x, \\
{}[x,y,z]&:=(xy)z-x(yz)=-[y,x,z]=-[x,z,y], \\
z\big(x(zy)\big)&=\big((zx)z\big)y,\qquad x\big(z(yz)\big)=\big((xz)y\big)z,\qquad (zx)(yz)=\big(z(xy)\big)z, \\
x\bar x&=\bar x x=N(x)\,e_0,\qquad \overline{xy}=\bar y\,\bar x,\qquad \mathrm{Re}(x)=\tfrac12(x+\bar x)=x_0, \\
N(xy)&=N(x)\,N(y),\qquad x\,x^{-1}=x^{-1}x=e_0\ \ \text{for}\ \ x^{-1}=\bar x/N(x), \\
x^2&-2\,\mathrm{Re}(x)\,x+N(x)\,e_0=0 .
\end{align*}
The last identity shows that the subalgebra generated by one octonion is $\mathbb R$ or $\mathbb C$ (power-associativity). Multiplicativity of the norm is Degen's eight-square identity: with $x=\sum a_Ae_A$ and $y=\sum b_Ae_A$ the components of $xy$ are
\begin{align*}
(xy)_{0}&=a_{0} b_{0} - a_{1} b_{1} - a_{2} b_{2} - a_{3} b_{3} - a_{4} b_{4} - a_{5} b_{5} - a_{6} b_{6} - a_{7} b_{7}, \\
(xy)_{1}&=a_{0} b_{1} + a_{1} b_{0} - a_{2} b_{3} + a_{3} b_{2} - a_{4} b_{5} + a_{5} b_{4} + a_{6} b_{7} - a_{7} b_{6}, \\
(xy)_{2}&=a_{0} b_{2} + a_{1} b_{3} + a_{2} b_{0} - a_{3} b_{1} - a_{4} b_{6} - a_{5} b_{7} + a_{6} b_{4} + a_{7} b_{5}, \\
(xy)_{3}&=a_{0} b_{3} - a_{1} b_{2} + a_{2} b_{1} + a_{3} b_{0} - a_{4} b_{7} + a_{5} b_{6} - a_{6} b_{5} + a_{7} b_{4}, \\
(xy)_{4}&=a_{0} b_{4} + a_{1} b_{5} + a_{2} b_{6} + a_{3} b_{7} + a_{4} b_{0} - a_{5} b_{1} - a_{6} b_{2} - a_{7} b_{3}, \\
(xy)_{5}&=a_{0} b_{5} - a_{1} b_{4} + a_{2} b_{7} - a_{3} b_{6} + a_{4} b_{1} + a_{5} b_{0} + a_{6} b_{3} - a_{7} b_{2}, \\
(xy)_{6}&=a_{0} b_{6} - a_{1} b_{7} - a_{2} b_{4} + a_{3} b_{5} + a_{4} b_{2} - a_{5} b_{3} + a_{6} b_{0} + a_{7} b_{1}, \\
(xy)_{7}&=a_{0} b_{7} + a_{1} b_{6} - a_{2} b_{5} - a_{3} b_{4} + a_{4} b_{3} + a_{5} b_{2} - a_{6} b_{1} + a_{7} b_{0},
\end{align*}
\begin{equation*}
\Big(\sum_{A=0}^{7}a_A^2\Big)\Big(\sum_{A=0}^{7}b_A^2\Big)=\sum_{L=0}^{7}(xy)_L^2 .
\end{equation*}
Hence $N(x)=0$ only for $x=0$, and $xy=0$ only if $x=0$ or $y=0$: no zero divisors, a division algebra (unlike the split octonions).

\textbf{Non-associativity.} Exactly 168 of the 512 ordered basis triples $(e_A,e_B,e_K)$ are non-associative: the triples of distinct imaginary units not on a common Fano line ($7\cdot6\cdot4=168$). Example: $(e_1e_2)e_4=e_7$ but $e_1(e_2e_4)=-e_7$.

\textbf{A numerical example.} For $x=e_{0}+2e_{1}-e_{4}+3e_{7}$, $y=-2e_{0}+e_{2}+e_{5}$, $z=e_{3}+e_{6}$:
\begin{align*}
xy&=-2e_{0}-3e_{1}+4e_{2}-2e_{3}+4e_{4}-2e_{5}-e_{6}-6e_{7}, \\
yx&=-2e_{0}-5e_{1}-2e_{2}+2e_{3}+4e_{5}+e_{6}-6e_{7}, \\
N(x)&=15,\quad N(y)=6,\quad N(xy)=90=N(x)N(y), \\
x^{-1}&=\frac{1}{15}e_{0}-\frac{2}{15}e_{1}+\frac{1}{15}e_{4}-\frac{1}{5}e_{7}, \\
(xy)z-x(yz)&=-6e_{1}-2e_{2}-6e_{3}+6e_{4}+2e_{5}+6e_{6}+6e_{7}\neq0 .
\end{align*}
\Needspace{8\baselineskip}\subsection*{6.5 Cayley–Dickson doubling \texorpdfstring{$\mathbb O=\mathbb H\oplus\ell\mathbb H$}{𝕆=ℍ⊕ℓℍ}}\addcontentsline{toc}{subsection}{6.5 Cayley–Dickson doubling \texorpdfstring{$\mathbb O=\mathbb H\oplus\ell\mathbb H$}{𝕆=ℍ⊕ℓℍ}}
The span of $e_0,e_1,e_2,e_3$ — the time direction and the three ordinary space directions $x_0,\dots,x_3$ — is a quaternion subalgebra $\mathbb H$, and $\ell=e_4$ doubles it: $e_4e_k=e_{4+k}$ ($k=1,2,3$), $e_4^2=-e_0$. Every octonion is $x=a+\ell b$ with quaternions $a=(x_0,x_1,x_2,x_3)$ and $b=(x_4,x_5,x_6,x_7)$. The orientation inherited from 4+4.nb is $e_1e_2=-e_3$ (so $i=e_1$, $j=e_2$, $k=-e_3$ are Hamilton's units). With $\ell^2=\gamma$ the product is
\begin{align*}
(a,b)(c,d)&=\big(ac+\gamma\,d\bar b,\ \ \bar a d+cb\big), \\
\gamma&=-1:\ \text{octonions (this document)},\qquad \gamma=+1:\ \text{split octonions (4+4.nb)}.
\end{align*}
Both cases are proved symbolically in 7+1.nb. The formula $(ac-\bar d b,\ \bar a d+cb)$ of [1A] fails: for $x=(1,1)$, $y=(i,0)$ it gives $x^2y=(0,2i)$ but $x(xy)=(2i,2i)$, and its norm is not multiplicative.

\Needspace{8\baselineskip}\subsection*{6.6 Octonions versus split octonions}\addcontentsline{toc}{subsection}{6.6 Octonions versus split octonions}
With the split units relabelled 0-based ($\varepsilon_{B+1}\to e^{S}_B$), the two tables differ only in the products of two doubled units:
\begin{equation*}
e_Ae_B\big|_{\mathbb O}=s_{AB}\;e^{S}_Ae^{S}_B\big|_{\text{split}},\qquad s_{AB}=\begin{cases}-1,&A\ge4\ \text{and}\ B\ge4,\\ +1,&\text{otherwise},\end{cases}
\end{equation*}
because $\ell^2=-1$ instead of $+1$. Over $\mathbb C$ the algebras coincide: $\varphi(e^{S}_A)=c_Ae_A$ with $c=(1,1,1,1,i,i,i,i)$ satisfies $\varphi(xy)=\varphi(x)\varphi(y)$ and $N_{\mathbb O}(\varphi(x))=N_S(x)$, so $\mathbb O_{\text{split}}\otimes\mathbb C\cong\mathbb O\otimes\mathbb C$. The real substitute for the factors $i$ is the matrix $\sigma$ of section 6.1.

\Needspace{12\baselineskip}\section*{7. Clifford algebras from the octonions}\addcontentsline{toc}{section}{7. Clifford algebras from the octonions}
\Needspace{8\baselineskip}\subsection*{7.1 The exact mirror of \texorpdfstring{$T_A$}{TA}: the Euclidean algebra \texorpdfstring{$Cl(8,0)$}{Cl(8,0)}}\addcontentsline{toc}{subsection}{7.1 The exact mirror of \texorpdfstring{$T_A$}{TA}: the Euclidean algebra \texorpdfstring{$Cl(8,0)$}{Cl(8,0)}}
Copying $T_A=\left(\begin{smallmatrix}0&\bar\tau_A\\\tau_A&0\end{smallmatrix}\right)$ with $\tau\to\tau^{O}$:
\begin{align*}
T^{O}_A&=\begin{pmatrix}0&\bar\tau^{O}_A\\ \tau^{O}_A&0\end{pmatrix},\qquad T^{O}_AT^{O}_B+T^{O}_BT^{O}_A=2\,g^{O}_{AB}\,1_{16}=2\,\delta_{AB}\,1_{16}\qquad(Cl(8,0)\cong\mathbb R(16)), \\
T^{O}_A&=T_A\ (A=0,1,2,3),\qquad T^{O}_k=T_k\,S_{44}\ (k=4,\dots,7),\qquad S_{44}=T_4T_5T_6T_7=\mathrm{diag}(\sigma,\sigma), \\
T^{O}_8&=T^{O}_0T^{O}_1\cdots T^{O}_7=T_8=\mathrm{diag}(-1_8,1_8).
\end{align*}
\Needspace{8\baselineskip}\subsection*{7.2 The real Clifford algebra of 7+1 spacetime: \texorpdfstring{$\Gamma^{(1,7)}$}{Γ(1,7)}, \texorpdfstring{$Cl(1,7)\cong\mathbb R(16)$}{Cl(1,7)≅ℝ(16)}}\addcontentsline{toc}{subsection}{7.2 The real Clifford algebra of 7+1 spacetime: \texorpdfstring{$\Gamma^{(1,7)}$}{Γ(1,7)}, \texorpdfstring{$Cl(1,7)\cong\mathbb R(16)$}{Cl(1,7)≅ℝ(16)}}
\textbf{Definition.}
\begin{align*}
\Gamma^{(1,7)}_0&=T^{O}_0=\begin{pmatrix}0&1_8\\1_8&0\end{pmatrix}\ \ \text{(time)},\qquad \Gamma^{(1,7)}_i=T^{O}_0\,T^{O}_i=\begin{pmatrix}\tau^{O}_i&0\\0&\bar\tau^{O}_i\end{pmatrix}\ \ (i=1,\dots,7,\ \text{space}), \\
\Gamma^{(1,7)}_0&=T_0,\qquad \Gamma^{(1,7)}_i=T_0\,T_i\ (i=1,2,3),\qquad \Gamma^{(1,7)}_k=T_0\,T_k\,S_{44}\ (k=4,\dots,7)\qquad\text{(only 4+4.nb matrices)}.
\end{align*}
\textbf{Relations} (all verified):
\begin{align*}
\Gamma^{(1,7)}_A\Gamma^{(1,7)}_B+\Gamma^{(1,7)}_B\Gamma^{(1,7)}_A&=2\,\eta^{(1,7)}_{AB}\,1_{16}=-2\,\eta^{(7,1)}_{AB}\,1_{16}, \\
(\Gamma^{(1,7)}_0)^{\mathsf T}&=\Gamma^{(1,7)}_0,\quad (\Gamma^{(1,7)}_i)^{\mathsf T}=-\Gamma^{(1,7)}_i, \\
(\Gamma^{(1,7)}_A)^{\mathsf T}&=\Gamma^{(1,7)}_0\,\Gamma^{(1,7)}_A\,\Gamma^{(1,7)}_0\quad\text{(real Dirac adjoint)}, \\
\mathrm{tr}\big(\Gamma^{(1,7)}_{A_1}\cdots\Gamma^{(1,7)}_{A_k}\big)&=0\quad(1\le k\le8,\ A_1<\dots<A_k).
\end{align*}
All entries are $0$ or $\pm1$. \textbf{Why this is the Clifford algebra of [1B].} (i) For $v=\sum_Av_Ae_A$ the matrix $\Gamma^{(1,7)}(v)=\sum_Av_A\Gamma^{(1,7)}_A$ satisfies
\begin{equation*}
\Gamma^{(1,7)}(v)^2=Q(v)\,1_{16},\qquad Q(v)=v^{\mathsf T}\eta^{(1,7)}v=v_0^2-v_1^2-\dots-v_7^2,
\end{equation*}
so by the universal property of $Cl(V,Q)=\mathcal T(V)/\mathcal I_Q(V)$ the map extends to an algebra homomorphism $Cl(V,Q)\to\mathbb R(16)$. (ii) The 256 ordered products $\Gamma^{(1,7)}_{A_1}\cdots\Gamma^{(1,7)}_{A_k}$ ($A_1<\dots<A_k$, $k=0,\dots,8$, counted by $\binom8k$) are linearly independent, so the homomorphism is injective; both sides have dimension 256, hence $Cl(1,7)\cong\mathbb R(16)$. (iii) Only multiples of $1_{16}$ commute with all $\Gamma^{(1,7)}_A$: the representation is irreducible over $\mathbb R$.

\Needspace{8\baselineskip}\subsection*{7.3 The complexified realization in the exact block form of 4+4.nb: \texorpdfstring{$\Gamma^{(7,1)}$}{Γ(7,1)}}\addcontentsline{toc}{subsection}{7.3 The complexified realization in the exact block form of 4+4.nb: \texorpdfstring{$\Gamma^{(7,1)}$}{Γ(7,1)}}
For the sign convention of 4+4.nb, $\Gamma_A\Gamma_B+\Gamma_B\Gamma_A=2\eta^{(7,1)}_{AB}1_{16}$, complex numbers are unavoidable (section 7.4). The complexification needed is minimal: multiply the time block by $i$.
\begin{align*}
\tau^{C}_0&=\bar\tau^{C}_0=i\,1_8,\qquad \tau^{C}_k=\tau^{O}_k,\quad \bar\tau^{C}_k=\bar\tau^{O}_k\ (k=1,\dots,7),\qquad \Gamma^{(7,1)}_A=\begin{pmatrix}0&\bar\tau^{C}_A\\ \tau^{C}_A&0\end{pmatrix}, \\
\Gamma^{(7,1)}_0&=i\,T_0,\qquad \Gamma^{(7,1)}_k=T^{O}_k\ (k=1,\dots,7),\qquad \Gamma^{(7,1)}_A\Gamma^{(7,1)}_B+\Gamma^{(7,1)}_B\Gamma^{(7,1)}_A=2\,\eta^{(7,1)}_{AB}\,1_{16}, \\
\Gamma^{(7,1)}_A&=S\,\big(i\,\Gamma^{(1,7)}_A\big)\,S^{-1},\qquad S=\tfrac{1}{\sqrt2}\big(1_{16}-i\,\Gamma^{(1,7)}_0\big),\qquad SS^{\dagger}=1_{16}.
\end{align*}
This is literally the formula $T_A=\left(\begin{smallmatrix}0&\bar\tau_A\\\tau_A&0\end{smallmatrix}\right)$ of 4+4.nb, applied to the complexified 8×8 matrices; only the time component is complex.

\Needspace{8\baselineskip}\subsection*{7.4 Theorem: with \texorpdfstring{$\eta^{(7,1)}$}{η(7,1)} there is no real 16×16 realization}\addcontentsline{toc}{subsection}{7.4 Theorem: with \texorpdfstring{$\eta^{(7,1)}$}{η(7,1)} there is no real 16×16 realization}
\textbf{Theorem.} There are no real 16×16 matrices $R_0,\dots,R_7$ with $R_AR_B+R_BR_A=2\eta^{(7,1)}_{AB}1_{16}$. Reordering, renaming or changing the signs of the indices cannot help, because none of these changes the signature $(7,1)$ of this relation: seven generators square to $+1$ and one to $-1$, the opposite of the squares of the octonion basis.

\textbf{Proof.} $Cl_8(\mathbb C)\cong\mathbb C(16)$ has exactly one irreducible module, $\mathbb C^{16}$. If real $R_A$ existed, $R_A$ and $\Gamma^{(7,1)}_A$ would be two 16-dimensional modules, so $R_A=T\,\Gamma^{(7,1)}_A\,T^{-1}$ for some invertible complex $T$. Since $R_A$ is real, complex conjugation gives
\begin{equation*}
M\,\overline{\Gamma^{(7,1)}_A}=\Gamma^{(7,1)}_A\,M\quad(A=0,\dots,7),\qquad M=T^{-1}\,\overline{T},\qquad M\,\overline M=1_{16}.
\end{equation*}
By Schur's lemma all such $M$ form a one-dimensional space. Solving the linear equations (7.4 of the notebook) gives exactly that space, spanned by the volume element
\begin{equation*}
M_0=\Gamma^{(1,7)}_8=\Gamma^{(1,7)}_0\Gamma^{(1,7)}_1\cdots\Gamma^{(1,7)}_7=\begin{pmatrix}0&-1_8\\1_8&0\end{pmatrix},\qquad M_0\,\overline{M_0}=M_0^2=-1_{16}.
\end{equation*}
Hence $M=cM_0$ and $M\overline M=-|c|^2\,1_{16}\neq1_{16}$, a contradiction. $\square$ The sign of $M\overline M$ is a Frobenius–Schur-type indicator: $+1$ means the complex module has a real form (the algebra is $\mathbb R(16)$), $-1$ means it is quaternionic (the algebra is $\mathbb H(8)$, whose smallest real module is $\mathbb H^8=\mathbb R^{32}$). For $\Gamma^{(1,7)}$ the indicator is $+1$ (take $M=1_{16}$).

\Needspace{8\baselineskip}\subsection*{7.5 A real realization of \texorpdfstring{$Cl(7,1)$}{Cl(7,1)} needs 32×32 matrices}\addcontentsline{toc}{subsection}{7.5 A real realization of \texorpdfstring{$Cl(7,1)$}{Cl(7,1)} needs 32×32 matrices}
\begin{align*}
\widehat\Gamma^{(7,1)}_A&=\varepsilon\otimes\Gamma^{(1,7)}_A,\qquad \varepsilon=\begin{pmatrix}0&1\\-1&0\end{pmatrix},\qquad \widehat\Gamma^{(7,1)}_A\widehat\Gamma^{(7,1)}_B+\widehat\Gamma^{(7,1)}_B\widehat\Gamma^{(7,1)}_A=2\,\eta^{(7,1)}_{AB}\,1_{32}, \\
Q_i&=\varepsilon\otimes1_{16},\qquad Q_j=\sigma_1\otimes\Gamma^{(1,7)}_8,\qquad Q_k=\sigma_3\otimes\Gamma^{(1,7)}_8, \\
Q_i^2&=Q_j^2=Q_k^2=-1_{32},\qquad Q_iQ_j=Q_k=-Q_jQ_i,\qquad Q\,\widehat\Gamma^{(7,1)}_A=\widehat\Gamma^{(7,1)}_A\,Q\ \ (Q=Q_i,Q_j,Q_k).
\end{align*}
The matrices commuting with all $\widehat\Gamma^{(7,1)}_A$ form exactly the 4-dimensional span of $1_{32},Q_i,Q_j,Q_k$ (computed in the notebook): the quaternions. This is $Cl(7,1)\cong\mathbb H(8)$ acting on $\mathbb R^{32}=\mathbb H^8$.

\Needspace{8\baselineskip}\subsection*{7.6 Volume elements and chirality}\addcontentsline{toc}{subsection}{7.6 Volume elements and chirality}
\begin{align*}
T_8^2&=+1_{16}\ \Rightarrow\ P_{L,R}=\tfrac12(1_{16}\mp T_8)\ \text{real (4+4.nb)};\qquad (T^{O}_8)^2=+1_{16},\ T^{O}_8=T_8, \\
\big(\Gamma^{(1,7)}_8\big)^2&=-1_{16},\qquad \Gamma^{(1,7)}_8\Gamma^{(1,7)}_A=-\Gamma^{(1,7)}_A\Gamma^{(1,7)}_8,\qquad \Gamma^{(1,7)}_8\Sigma_{AB}=\Sigma_{AB}\Gamma^{(1,7)}_8, \\
\chi^{(7,1)}&=-i\,\Gamma^{(7,1)}_8=T^{O}_8=T_8,\qquad \tfrac12\big(1_{16}\mp\chi^{(7,1)}\big)=P_{L,R}, \\
\chi^{(1,7)}&=i\,\Gamma^{(1,7)}_8,\qquad \tfrac12\big(1_{16}\mp\chi^{(1,7)}\big)\ \text{complex projectors of rank 8 each}.
\end{align*}
With real matrices in 7+1 dimensions there is no real chirality projector: $\Gamma^{(1,7)}_8$ is a complex structure commuting with Spin(1,7), so the real 16-component spinor is one complex 8-component Weyl spinor. In the basis of $\Gamma^{(7,1)}$ the chiral projectors are literally $P_L$ and $P_R$ of 4+4.nb.

\Needspace{8\baselineskip}\subsection*{7.7 What the octonions do: \texorpdfstring{$\Gamma(x)^2=\mathrm{Re}(xx)$}{Γ(x)2=Re(xx)} and alternativity}\addcontentsline{toc}{subsection}{7.7 What the octonions do: \texorpdfstring{$\Gamma(x)^2=\mathrm{Re}(xx)$}{Γ(x)2=Re(xx)} and alternativity}
\begin{align*}
\mathrm{Re}(x\,x)&=x_0^2-x_1^2-\dots-x_7^2=x^{\mathsf T}\eta^{(1,7)}x,\qquad \Gamma^{(1,7)}(x)^2=\mathrm{Re}(x\,x)\,1_{16}, \\
T^{O}(x)^2&=N(x)\,1_{16}=x^{\mathsf T}g^{O}x\,1_{16},\qquad T(x)^2=x^{\mathsf T}\eta^{(4,4)}x\,1_{16}\ \text{(4+4.nb)}.
\end{align*}
Rotating both 8-blocks by $E^{O}$ gives $\Gamma^{(1,7)}$ in the octonion basis,
\begin{equation*}
\Gamma^{e}_0=\begin{pmatrix}0&1_8\\1_8&0\end{pmatrix},\quad \Gamma^{e}_i=\begin{pmatrix}L^{O}_i&0\\0&-L^{O}_i\end{pmatrix},\qquad \Gamma^{e}(x)\begin{pmatrix}\psi_1\\ \psi_2\end{pmatrix}=\begin{pmatrix}x_0\psi_2+\vec x\,\psi_1\\ x_0\psi_1-\vec x\,\psi_2\end{pmatrix}\quad(\psi_1,\psi_2\in\mathbb O),
\end{equation*}
\begin{equation*}
\Gamma^{e}(x)^2\begin{pmatrix}\psi_1\\ \psi_2\end{pmatrix}=\begin{pmatrix}x_0^2\psi_1+\vec x(\vec x\psi_1)\\ x_0^2\psi_2+\vec x(\vec x\psi_2)\end{pmatrix}=(x_0^2-|\vec x|^2)\begin{pmatrix}\psi_1\\ \psi_2\end{pmatrix},
\end{equation*}
where $\vec x=\sum_{i}x_ie_i$ and the last step is alternativity, $\vec x(\vec x\psi)=(\vec x\vec x)\psi=-|\vec x|^2\psi$. The Clifford relation of 7+1 spacetime is the alternativity of the octonions.

\Needspace{12\baselineskip}\section*{8. Examples that test everything}\addcontentsline{toc}{section}{8. Examples that test everything}
\Needspace{8\baselineskip}\subsection*{8.1 The Lorentz algebra so(1,7)}\addcontentsline{toc}{subsection}{8.1 The Lorentz algebra so(1,7)}
\begin{align*}
\Sigma_{AB}&=\tfrac14\big(\Gamma^{(1,7)}_A\Gamma^{(1,7)}_B-\Gamma^{(1,7)}_B\Gamma^{(1,7)}_A\big)\quad(A<B:\ 28\ \text{real, linearly independent matrices}), \\
\big[\Sigma_{AB},\Gamma^{(1,7)}_K\big]&=\eta^{(1,7)}_{BK}\,\Gamma^{(1,7)}_A-\eta^{(1,7)}_{AK}\,\Gamma^{(1,7)}_B, \\
\big[\Sigma_{AB},\Sigma_{KL}\big]&=\eta^{(1,7)}_{BK}\Sigma_{AL}-\eta^{(1,7)}_{AK}\Sigma_{BL}-\eta^{(1,7)}_{BL}\Sigma_{AK}+\eta^{(1,7)}_{AL}\Sigma_{BK}.
\end{align*}
\Needspace{8\baselineskip}\subsection*{8.2 A boost along \texorpdfstring{$x_1$}{x1} and a rotation in the \texorpdfstring{$x_4$}{x4}–\texorpdfstring{$x_5$}{x5} plane}\addcontentsline{toc}{subsection}{8.2 A boost along \texorpdfstring{$x_1$}{x1} and a rotation in the \texorpdfstring{$x_4$}{x4}–\texorpdfstring{$x_5$}{x5} plane}
Because $(\Gamma^{(1,7)}_0\Gamma^{(1,7)}_1)^2=+1_{16}$ and $(\Gamma^{(1,7)}_4\Gamma^{(1,7)}_5)^2=-1_{16}$,
\begin{align*}
S(\beta)&=e^{\beta\Sigma_{01}}=\cosh\tfrac\beta2\,1_{16}+\sinh\tfrac\beta2\,\Gamma^{(1,7)}_0\Gamma^{(1,7)}_1,\qquad S(\beta)^{-1}=\cosh\tfrac\beta2\,1_{16}-\sinh\tfrac\beta2\,\Gamma^{(1,7)}_0\Gamma^{(1,7)}_1, \\
R(\theta)&=e^{\theta\Sigma_{45}}=\cos\tfrac\theta2\,1_{16}+\sin\tfrac\theta2\,\Gamma^{(1,7)}_4\Gamma^{(1,7)}_5,\qquad R(\theta)^{-1}=\cos\tfrac\theta2\,1_{16}-\sin\tfrac\theta2\,\Gamma^{(1,7)}_4\Gamma^{(1,7)}_5,
\end{align*}
both real (they solve $dS/d\beta=\Sigma_{01}S$, $S(0)=1_{16}$ and $dR/d\theta=\Sigma_{45}R$, $R(0)=1_{16}$). They act on the generators as Lorentz transformations, with $\Lambda$ read off from traces:
\begin{equation*}
S\,\Gamma^{(1,7)}_K\,S^{-1}=\sum_{B=0}^{7}\Lambda_{BK}\,\Gamma^{(1,7)}_B,\qquad \Lambda_{BK}=\tfrac1{16}\,\mathrm{tr}\big(\eta^{(1,7)}_{BB}\,\Gamma^{(1,7)}_B\,S\,\Gamma^{(1,7)}_K\,S^{-1}\big),\qquad \Lambda^{\mathsf T}\eta^{(1,7)}\Lambda=\eta^{(1,7)},
\end{equation*}
\begin{equation*}
\Lambda_{\text{boost}}=\begin{pmatrix}\cosh\beta&-\sinh\beta&0\\-\sinh\beta&\cosh\beta&0\\0&0&1_6\end{pmatrix},\qquad \Lambda_{\text{rot}}=\begin{pmatrix}1_4&0&0\\0&\begin{smallmatrix}\cos\theta&-\sin\theta\\ \sin\theta&\cos\theta\end{smallmatrix}&0\\0&0&1_2\end{pmatrix}.
\end{equation*}
\Needspace{8\baselineskip}\subsection*{8.3 The 7+1 Dirac operator squares to the wave operator}\addcontentsline{toc}{subsection}{8.3 The 7+1 Dirac operator squares to the wave operator}
\begin{equation*}
\Gamma^{A}=\eta_{(1,7)}^{AA}\,\Gamma^{(1,7)}_A,\qquad \big(\Gamma^{A}\partial_A\big)^2\Psi=\eta_{(1,7)}^{AB}\,\partial_A\partial_B\Psi=\big(\partial_0^2-\partial_1^2-\dots-\partial_7^2\big)\Psi
\end{equation*}
for 16 completely arbitrary functions $\Psi=(\Psi_1,\dots,\Psi_{16})$ of $x_0,\dots,x_7$ (the notebook proves it symbolically; here $\eta_{(1,7)}^{AB}$ is the inverse metric, equal to $\eta^{(1,7)}_{AB}$).

\Needspace{8\baselineskip}\subsection*{8.4 Spinors on the mass shell: real amplitudes}\addcontentsline{toc}{subsection}{8.4 Spinors on the mass shell: real amplitudes}
Take $p=(3,2,2,0,0,0,0,0)$ (contravariant components $p^A$) and $m=1$: $p^{\mathsf T}\eta^{(1,7)}p=9-4-4=1=m^2$. With $\not p=\sum_Ap^A\Gamma^{(1,7)}_A$ we have $\not p^{\,2}=m^2\,1_{16}$, so $\not p$ has eigenvalues $\pm m$, eight times each. Because $\not p$ is a real matrix, $(\not p-m)w=0$ has an 8-dimensional space of REAL solutions $w$ (spinor amplitudes). The basis printed by Mathematica (\texttt{NullSpace}, one solution per row):
\begin{equation*}
W=\begin{pmatrix}2 & 2 & 0 & 0 & 0 & 0 & 0 & 1 & 0 & 0 & 0 & 0 & 0 & 0 & 0 & 3 \\ -2 & 2 & 0 & 0 & 0 & 0 & 1 & 0 & 0 & 0 & 0 & 0 & 0 & 0 & 3 & 0 \\ 0 & 0 & -2 & -2 & 0 & 1 & 0 & 0 & 0 & 0 & 0 & 0 & 0 & 3 & 0 & 0 \\ 0 & 0 & 2 & -2 & 1 & 0 & 0 & 0 & 0 & 0 & 0 & 0 & 3 & 0 & 0 & 0 \\ 0 & 0 & 0 & 1 & 2 & 2 & 0 & 0 & 0 & 0 & 0 & 3 & 0 & 0 & 0 & 0 \\ 0 & 0 & 1 & 0 & -2 & 2 & 0 & 0 & 0 & 0 & 3 & 0 & 0 & 0 & 0 & 0 \\ 0 & 1 & 0 & 0 & 0 & 0 & -2 & -2 & 0 & 3 & 0 & 0 & 0 & 0 & 0 & 0 \\ 1 & 0 & 0 & 0 & 0 & 0 & 2 & -2 & 3 & 0 & 0 & 0 & 0 & 0 & 0 & 0\end{pmatrix},\qquad (\not p-m)\,W^{\mathsf T}=0 .
\end{equation*}
Each row $w$ gives the plane wave $\Psi=w\,e^{-i\,p\cdot x}$, $p\cdot x=p^{\mathsf T}\eta^{(1,7)}x$, solving $(i\Gamma^{A}\partial_A-m)\Psi=0$. The plane wave itself is complex: real amplitudes do not make a real field (section 8.5).

\Needspace{8\baselineskip}\subsection*{8.5 Real spinor fields: massless with 16 components, massive with 32}\addcontentsline{toc}{subsection}{8.5 Real spinor fields: massless with 16 components, massive with 32}
The operator $i\Gamma^{A}\partial_A$ is imaginary, so a real $\Psi$ cannot solve $(i\Gamma^{A}\partial_A-m)\Psi=0$. A real first-order equation $(G^{A}\partial_A-m)\Psi=0$ with real matrices $G_A$ ($G^{A}=\eta_{(7,1)}^{AA}G_A$) implies the Klein–Gordon equation, through
\begin{equation*}
\big(G^{A}\partial_A+m\big)\big(G^{A}\partial_A-m\big)=-\big(\partial_0^2-\partial_1^2-\dots-\partial_7^2+m^2\big),
\end{equation*}
only if $G_AG_B+G_BG_A=2\,\eta^{(7,1)}_{AB}\,1$ (time squaring to $-1$), and that relation has no real 16×16 solution (section 7.4). The remaining real option, $(\Gamma^{A}\partial_A-mJ)\Psi=0$, squares to the Klein–Gordon operator $\partial_0^2-\partial_1^2-\dots-\partial_7^2+m^2$ only if the real matrix $J$ anticommutes with every $\Gamma^{(1,7)}_A$ and $J^2=+1_{16}$. But (computed in the notebook)
\begin{equation*}
\big\lbrace J\in\mathbb R(16):\ J\,\Gamma^{(1,7)}_A=-\Gamma^{(1,7)}_A\,J\ \ (A=0,\dots,7)\big\rbrace=\mathbb R\,\Gamma^{(1,7)}_8,\qquad \big(c\,\Gamma^{(1,7)}_8\big)^2=-c^2\,1_{16},
\end{equation*}
and with $J=c\,\Gamma^{(1,7)}_8$ the square is $\partial_0^2-\partial_1^2-\dots-\partial_7^2-c^2m^2$, a tachyon. So with 16 real components only massless real fields exist. For the null momentum $k=(1,1,0,0,0,0,0,0)$, $k^{\mathsf T}\eta^{(1,7)}k=0$ and $\not k=\sum_Ak^A\Gamma^{(1,7)}_A$ satisfies $\not k^{\,2}=0$; the real equation $\not k\,w=0$ has 8 real solutions, and for each of them
\begin{equation*}
\Psi=w\cos(k\cdot x),\qquad \Gamma^{A}\partial_A\Psi=-\sin(k\cdot x)\,\not k\,w=0\qquad(\text{a real solution}).
\end{equation*}
A real massive field needs the real 32×32 matrices $\widehat\Gamma^{(7,1)}_A=\varepsilon\otimes\Gamma^{(1,7)}_A$ of section 7.5, which do satisfy the $\eta^{(7,1)}$ relation. With $\widehat\Gamma^{A}=\eta_{(7,1)}^{AA}\widehat\Gamma^{(7,1)}_A$ and $P=\sum_Ap^A\widehat\Gamma^{(7,1)}_A$ ($p$ as in section 8.4):
\begin{align*}
\big(\widehat\Gamma^{A}\partial_A\big)^2&=-\big(\partial_0^2-\partial_1^2-\dots-\partial_7^2\big)\,1_{32},\qquad P^2=\big(p^{\mathsf T}\eta^{(7,1)}p\big)\,1_{32}=-m^2\,1_{32}, \\
\Psi&=w\cos(p\cdot x)+\tfrac1m\,P\,w\,\sin(p\cdot x)\qquad(w\in\mathbb R^{32}\ \text{arbitrary}), \\
\big(\widehat\Gamma^{A}\partial_A-m\big)\Psi&=0,\qquad \big(\partial_0^2-\partial_1^2-\dots-\partial_7^2+m^2\big)\Psi=0.
\end{align*}
Derivation: with $\phi=p\cdot x$ one has $\partial_A\phi=(\eta^{(1,7)}p)_A$ and $\eta_{(7,1)}^{AA}\eta^{(1,7)}_{AA}=-1$, so $\widehat\Gamma^{A}\partial_A\Psi=-P\big(-w\sin\phi+\tfrac1mPw\cos\phi\big)=Pw\sin\phi+m\,w\cos\phi=m\Psi$. The 32 real components are the quaternionic module $\mathbb H^8=\mathbb R^{32}$ of $Cl(7,1)\cong\mathbb H(8)$.

\Needspace{8\baselineskip}\subsection*{8.6 All nine signatures \texorpdfstring{$p+q=8$}{p+q=8}}\addcontentsline{toc}{subsection}{8.6 All nine signatures \texorpdfstring{$p+q=8$}{p+q=8}}
Complex generators of every signature: the first $p$ of the $T^{O}_A$ and $i$ times the remaining $q$. The indicator of section 7.4 then decides $\mathbb R(16)$ or $\mathbb H(8)$ (table of section 4.3). Real 16×16 generators exist exactly in the five $\mathbb R(16)$ cases, and all five are built here from the 4+4 and octonion matrices:
\begin{align*}
Cl(8,0)&:\ T^{O}_A,\qquad Cl(4,4):\ T_A\ \text{(4+4.nb)},\qquad Cl(1,7):\ \Gamma^{(1,7)}_A, \\
Cl(5,3)&:\ \Gamma^{(5,3)}_0=T_0,\ \ \Gamma^{(5,3)}_i=T_0T_i\ (i=1,\dots,7),\qquad \eta^{(5,3)}=\mathrm{diag}(1,-1,-1,-1,1,1,1,1), \\
&\phantom{:}\ \Gamma^{(5,3)}_A\Gamma^{(5,3)}_B+\Gamma^{(5,3)}_B\Gamma^{(5,3)}_A=2\,\eta^{(5,3)}_{AB}1_{16}, \\
Cl(0,8)&:\ \Gamma^{(0,8)}_A=T^{O}_8\,T^{O}_A,\qquad \Gamma^{(0,8)}_A\Gamma^{(0,8)}_B+\Gamma^{(0,8)}_B\Gamma^{(0,8)}_A=-2\,\delta_{AB}1_{16}.
\end{align*}
$\Gamma^{(5,3)}$ is the split-octonion analogue of $\Gamma^{(1,7)}$: its quadratic form $x^{\mathsf T}\eta^{(5,3)}x$ is $\mathrm{Re}(xx)$ of the split octonions, just as $x^{\mathsf T}\eta^{(1,7)}x$ is $\mathrm{Re}(xx)$ of the octonions.

\Needspace{12\baselineskip}\section*{9. Verification}\addcontentsline{toc}{section}{9. Verification}
\textbf{In Mathematica.} 7+1.nb was evaluated in 15.0.1 for Microsoft Windows (64-bit) (July 2, 2026) by \texttt{tools/run\_eval.wl} (all cells evaluated in a fresh kernel, then saved), in the private context \texttt{Notebook\$\$59\$143224}. It contains 166 input cells and 399 checks \texttt{check["name", test]}: 363 print their result where they stand, and 36 are made by the function \texttt{frameChecks} of section 13.1 (nine checks for each of four metrics) and printed by the summary table of section 24. All 399 returned True, the summary cells printed \texttt{failedChecks = \{\}} and \texttt{allChecksPassed = True}, there were 0 message cells, and the evaluation took 271.569 s.

{\small
\begin{longtable}{>{\raggedright\arraybackslash}p{\dimexpr 0.7333\linewidth-2\tabcolsep\relax}>{\raggedright\arraybackslash}p{\dimexpr 0.1333\linewidth-2\tabcolsep\relax}>{\raggedright\arraybackslash}p{\dimexpr 0.1333\linewidth-2\tabcolsep\relax}}
\toprule
\textbf{notebook section} & \textbf{checks} & \textbf{True} \\
\midrule\endfirsthead
\toprule
\textbf{notebook section} & \textbf{checks} & \textbf{True} \\
\midrule\endhead
\endfoot
\bottomrule\endlastfoot
2. Recap: the 4+4 objects, reproduced exactly from 4+4.nb & 24 & 24 \\
3. The octonions: a real realization mirroring the split octonions & 47 & 47 \\
4. Clifford algebras from the octonions: 16×16 matrices & 33 & 33 \\
5. Examples that test everything & 25 & 25 \\
7. The Dirac equation of the quantized 16-component field dirac16complex & 35 & 35 \\
8. A 16-component Kohn–Sham theory of the self-coupled field dirac16complex & 23 & 23 \\
9. The chiral Kohn–Sham theory (8 + 8 components) & 16 & 16 \\
11. Gravity on the 7+1 spacetime: metric, curvature and the Einstein tensor & 7 & 7 \\
12. The Lovelock tensors of 7+1 spacetime & 28 & 28 \\
13. Generalized frame matrices, the canonical spin connection, and the covariant Dirac equation of dirac16complex & 51 & 51 \\
14. The three epochs: primordial (near x0 = 0), modern, far future & 14 & 14 \\
15. The generalized self-coupling of dirac16complex: all nine ranks & 16 & 16 \\
16. The Kohn–Sham theory of dirac16complex in gravitational fields & 8 & 8 \\
17. The Standard Model on the world brane of 7+1 spacetime & 11 & 11 \\
18. dirac16complex at the world brane: the 16 components as four brane Dirac spinors, and the portals & 7 & 7 \\
19. Neutrino flavour and flavour conversion & 6 & 6 \\
20. The coupled field equations: every interaction & 14 & 14 \\
21. Regime I: the primordial Lovelock regime (near x0 = 0) & 19 & 19 \\
22. Regime II: the modern regime (a stabilized auxiliary space, four-dimensional gravity on the brane) & 9 & 9 \\
23. Regime III: the far future & 6 & 6 \\*
\end{longtable}}
\textbf{Independently in Python} (exact arithmetic with sympy, in \texttt{tools/}): \texttt{ref44.py} rebuilds every 4+4.nb object and reproduces the split table of 4+4.nb; \texttt{ref71.py} and \texttt{extra71.py} rebuild every 7+1 object and re-prove the relations; \texttt{verify\_7p1\_nb.py} parses the evaluated 7+1.nb and confirms that the matrices $\tau_A$, $\tau^{O}_A$, $\Gamma^{(1,7)}_A$, the two multiplication tables, the matrices $\Lambda$, the eight real amplitudes $W$ and the classification table printed by Mathematica coincide with the Python values. For sections 7–9, \texttt{ref\_dirac16.py} repeats every construction (sympy, mpmath with 30 digits, numpy) with 90 checks, all True, and \texttt{verify\_7p1\_nb.py} compares the printed spinors, the printed closed-form integrals and 23 printed numbers with it (agreement to at least 17 digits) and confirms that sections 0–6 are cell for cell the notebook of the first stage, except the cells of the explanation of the signature added on 27 October 2026 (section 2).

For sections 11–19, the independent programs \texttt{tools/ref\_gravity.py} (101 checks, 101 True; with its own curvature and Lovelock engine \texttt{tools/lovelock71.py}), \texttt{tools/ref\_ks\_gravity.py} (29 checks, 29 True) and \texttt{tools/ref\_sm.py} (33 checks, 33 True) repeat every construction; the Lovelock tensors are computed in the notebook both with xAct (as in the reference notebook Einstein-Lovelock-7+1-Nash.nb) and with an independent pair-set sum, which agree exactly.

For sections 20–23 the independent program \texttt{tools/ref\_coupled.py} (49 checks, 49 True; numpy, scipy, sympy, mpmath, with its own closed forms from \texttt{tools/lovelock71.py}, its own quadratures and an 8th-order Runge–Kutta integrator) repeats every derivation and every numerical solution; the 34 groups of numbers printed by both programs agree to a largest relative difference of 0.00032, and \texttt{verify\_7p1\_nb.py} confirms that every cell and output of sections 0–19 is unchanged.

\Needspace{12\baselineskip}\section*{10. The notebook 7+1.nb}\addcontentsline{toc}{section}{10. The notebook 7+1.nb}
\begin{itemize}
\item \textbf{0. How to use this notebook} — evaluation, the \texttt{check} bookkeeping, the private context, index conventions.
\item \textbf{1. Conventions, coordinates and notation} — the coordinate table, the metrics $\eta^{(4,4)}$, $\eta^{(7,1)}$, $\eta^{(1,7)}$, $g^{O}$, the dictionary of names, the definitions [1A], [1B] and the four corrections.
\item \textbf{2. Recap of 4+4.nb} — $\sigma$, $s_h$, $t_h$, $\zeta_A$, $\tau_A$, $\bar\tau_A$, $T_A$, $P_L$, $P_R$, $u$, $E_{44}$, the split table \texttt{gridx44} compared with the table printed by 4+4.nb, split product, zero divisors.
\item \textbf{3. The octonions} — $\tau^{O}$, vielbein $E^{O}$, \texttt{gridO}, Fano plane (with a picture), the laws of [1A], the eight-square identity, non-associativity, a numerical example, Cayley–Dickson, comparison with the split octonions.
\item \textbf{4. Clifford algebras} — $T^{O}$ ($Cl(8,0)$), $\Gamma^{(1,7)}$ (real 7+1), $\Gamma^{(7,1)}$ (complexified 7+1), the theorem, $\widehat\Gamma^{(7,1)}$ (real 32×32) and its quaternionic commutant, chirality, the octonionic meaning.
\item \textbf{5. Examples} — so(1,7), boost and rotation, Dirac operator, real on-shell spinor amplitudes, real massless (16 components) and massive (32 components) spinor fields, all nine signatures.
\item \textbf{6. Summary of all checks} of sections 1–5 — \texttt{Length[checks]}, \texttt{failedChecks}, \texttt{allChecksPassed} and a table of the checks.
\item \textbf{7. The Dirac equation of the quantized 16-component field dirac16complex} — real Dirac matrices, Dirac and Klein–Gordon equations, Hamiltonian, real spinors, spin sums, Grassmann algebra, quantization, charge conjugation, Lorentz covariance, chirality and the two coupled 8-component equations (section 11.1–11.2).
\item \textbf{8. A 16-component Kohn–Sham theory of the self-coupled field} — Lagrangian, nonlinear field equation, Hohenberg–Kohn, Kohn–Sham equations, exchange (Wick), Fierz, uniform matter (section 11.3, 11.5).
\item \textbf{9. The chiral Kohn–Sham theory (8 + 8 components)} — chiral densities, chiral Kohn–Sham equations, chiral symmetry, the chiral vacuum with the Dirac sea (section 11.4–11.5).
\item \textbf{10. Summary of all checks} of sections 1–9.
\item \textbf{11. Gravity on the 7+1 spacetime} — the curvature engine, the cosmological metric, symmetries and the Bianchi identity (section 12.1).
\item \textbf{12. The Lovelock tensors} — Lovelock’s theorem, the generalized Kronecker delta, xAct (as in Einstein-Lovelock-7+1-Nash.nb), the pair-set algorithm, closed forms, an inhomogeneous metric, a gravitational wave, exponential vacuum solutions (section 12.1).
\item \textbf{13. Generalized frame matrices, the canonical spin connection and the covariant Dirac equation} — total covariant derivative, spin curvature, Lichnerowicz, Cartan’s structure equation, a static radial metric, gauge coupling (section 12.2).
\item \textbf{14. The three epochs} — Kasner solutions near $x_0=0$, pair creation, a static auxiliary space, de Sitter and thermal pair creation (section 12.3).
\item \textbf{15. The generalized self-coupling} — nine ranks, the Fierz matrix, the covariant nonlinear field equation, chirality (section 12.4).
\item \textbf{16. Kohn–Sham theory in gravitational fields} — uniform matter in the cosmological spacetime, the epochs, self-consistent Einstein–Lovelock evolution (section 12.4).
\item \textbf{17. The Standard Model on the world brane} — fields, anomalies, Lagrangian, electroweak breaking, masses and mixing, cosmology (section 12.5).
\item \textbf{18. dirac16complex at the brane} — four brane Dirac spinors, Kaluza–Klein masses, portals, coupled field equations (section 12.6).
\item \textbf{19. Neutrino flavour} — vacuum, matter, the Sun, supernovae, the relic background, effects of dirac16complex (section 12.6).
\item \textbf{20. The coupled field equations: every interaction} — the action term by term (the cubic Lovelock density in invariant form), the energy–momentum tensors, the homogeneous equations, the radion and the four-dimensional constants (section 13.1).
\item \textbf{21. Regime I: the primordial Lovelock regime} — the exact vacuum curve, the cubic, Gauss–Bonnet and Einstein solutions and their accuracy, the matter-dominated cubic regime, decompactification, pair creation (section 13.2).
\item \textbf{22. Regime II: the modern regime} — capture of the radion, its decay and the moduli problem, $g_*(T)$ and $t(T)$, the late universe and its age (section 13.3).
\item \textbf{23. Regime III: the far future} — de Sitter, the radion in de Sitter, the size of every interaction (section 13.4).
\item \textbf{24. Summary of all checks} of sections 1–23.
\end{itemize}
To re-evaluate: open 7+1.nb in Wolfram 15.0.1, start from a fresh kernel (Evaluation \texorpdfstring{\ensuremath{\triangleright}}{▸} Quit Kernel \texorpdfstring{\ensuremath{\triangleright}}{▸} Local, if the kernel has been used: the tensor packages xAct of section 12 keep their definitions in the kernel) and choose Evaluation \texorpdfstring{\ensuremath{\triangleright}}{▸} Evaluate Notebook (about 272 seconds); or run \texttt{Get["D:/Developer/github/7+1/tools/run\_eval.wl"]} from any notebook, which evaluates 7+1.nb in place in a fresh kernel (a front-end kernel configuration "Eval71" that it adds for the purpose and removes again) and saves it. The notebook is generated by \texttt{tools/make\_7p1\_nb.py} (\texttt{python3 tools/make\_7p1\_nb.py} writes the unevaluated 7+1.nb).

\Needspace{12\baselineskip}\section*{11. Sections 7–10: the Dirac equation and the Kohn–Sham theories of dirac16complex}\addcontentsline{toc}{section}{11. Sections 7–10: the Dirac equation and the Kohn–Sham theories of dirac16complex}
Sections 7–10 of 7+1.nb were appended in a second stage (sections 0–6 are unchanged, except for the explanation of the signature of section 2, added on 27 October 2026). They contain the Dirac equation of a 16-component quantized fermion field named dirac16complex on the real 7+1 spacetime, with eight REAL 16×16 Dirac matrices; its chiral form as two 8-component equations; and the 16-component and chiral Kohn–Sham theories of that field coupled to itself. This section states the results with the notation of this document (section 3); 74 checks of 7+1.nb verify them.

\Needspace{8\baselineskip}\subsection*{11.1 The Dirac equation with real Dirac matrices}\addcontentsline{toc}{subsection}{11.1 The Dirac equation with real Dirac matrices}
The field is the column $\Psi=(\Psi_1,\dots,\Psi_{16})^{\mathsf T}$ (\texttt{dirac16complex}) of complex anticommuting (Grassmann-odd) operators, with the equal-time anticommutators
\begin{equation*}
\Psi_a(x_0,\vec x)\Psi_b(x_0,\vec y)+\Psi_b(x_0,\vec y)\Psi_a(x_0,\vec x)=0,\qquad \Psi_a(x_0,\vec x)\Psi_b^\dagger(x_0,\vec y)+\Psi_b^\dagger(x_0,\vec y)\Psi_a(x_0,\vec x)=\delta_{ab}\,\delta^7(\vec x-\vec y),
\end{equation*}
where $\vec x=(x_1,\dots,x_7)$ here denotes the space coordinates. The Dirac matrices are the real generators of section 7.2, $\Gamma^{(1,7)}_0=\left(\begin{smallmatrix}0&1_8\\1_8&0\end{smallmatrix}\right)$, $\Gamma^{(1,7)}_i=\left(\begin{smallmatrix}\tau^{O}_i&0\\0&-\tau^{O}_i\end{smallmatrix}\right)$, with $\Gamma^{A}=\eta_{(1,7)}^{AA}\Gamma^{(1,7)}_A$ and $\bar\Psi=\Psi^\dagger\Gamma^{0}$:
\begin{align*}
&\big(i\,\Gamma^{A}\partial_A-m\big)\Psi=0,\qquad \big(i\Gamma^{A}\partial_A+m\big)\big(i\Gamma^{A}\partial_A-m\big)=-\big(\partial_0^2-\partial_1^2-\dots-\partial_7^2+m^2\big)1_{16}, \\
&i\,\partial_0\Psi=H_{\rm D}\Psi,\qquad H_{\rm D}=-i\,\Gamma^{0}\Gamma^{i}\partial_i+m\,\Gamma^{0},\qquad \partial_A\big(\bar\Psi\Gamma^{A}\Psi\big)=0 .
\end{align*}
$\Gamma^0$ and $\Gamma^0\Gamma^i$ are real symmetric, so $H_{\rm D}$ is Hermitian. The factor $i$ is required: with real matrices, $(\Gamma^A\partial_A)^2=\partial_0^2-\partial_1^2-\dots-\partial_7^2$, and only $i\Gamma^A\partial_A$ gives the Klein–Gordon operator with $+m^2$. For a plane wave $e^{-i(Ex_0-\vec p\cdot\vec x)}$, $H_{\rm D}$ becomes the real matrix $\left(\begin{smallmatrix}0&m1_8+p^i\tau^{O}_i\\m1_8-p^i\tau^{O}_i&0\end{smallmatrix}\right)$, and since $(m1_8+p^i\tau^{O}_i)(m1_8-p^i\tau^{O}_i)=E^21_8$, $E=\sqrt{m^2+\vec p^{\,2}}$, its positive-energy eigenvectors are REAL: $w_s=\frac1{\sqrt E}\big(E\,e_s;\ (m1_8-p^i\tau^{O}_i)e_s\big)$, $s=1,\dots,8$ ($e_s$ the unit columns of length 8), real amplitudes with $(\not p-m)w_s=0$ as in section 8.4. The notebook also verifies the spin sums, the mode expansion with the canonical anticommutators, charge conjugation $\Psi^c=\Gamma^{(1,7)}_8\Psi^*$ and Lorentz covariance.

\Needspace{8\baselineskip}\subsection*{11.2 Chiral projections: two coupled 8-component equations}\addcontentsline{toc}{subsection}{11.2 Chiral projections: two coupled 8-component equations}
There is no real chirality matrix (section 8.5: the real matrices anticommuting with all $\Gamma^{(1,7)}_A$ are the multiples of $\Gamma^{(1,7)}_8$, which squares to $-1_{16}$). With the complex chirality $\chi^{(1,7)}=i\Gamma^{(1,7)}_8$ (section 7.6) the chiral projections $P^{(1,7)}_{L,R}=\frac12\big(1_{16}\mp\chi^{(1,7)}\big)$ give, with $\Psi=(\Psi_\uparrow;\Psi_\downarrow)$ (upper and lower 8 components),
\begin{align*}
\Psi_L&=\frac{\Psi_\uparrow+i\,\Psi_\downarrow}{\sqrt2},\qquad \Psi_R=\frac{i\,\Psi_\uparrow+\Psi_\downarrow}{\sqrt2}, \\
P^{(1,7)}_L\Psi&=\frac1{\sqrt2}\begin{pmatrix}1_8\\-i1_8\end{pmatrix}\Psi_L,\qquad P^{(1,7)}_R\Psi=\frac1{\sqrt2}\begin{pmatrix}-i1_8\\1_8\end{pmatrix}\Psi_R, \\
\big(i\,\partial_0+\tau^{O}_i\,\partial_i\big)\Psi_L&=m\,\Psi_R,\qquad \big(i\,\partial_0-\tau^{O}_i\,\partial_i\big)\Psi_R=m\,\Psi_L .
\end{align*}
The two 8-component equations contain only the real matrices $\tau^{O}_i$; the mass couples them, and for $m=0$ they decouple.

\Needspace{8\baselineskip}\subsection*{11.3 The self-coupled field and its 16-component Kohn–Sham theory}\addcontentsline{toc}{subsection}{11.3 The self-coupled field and its 16-component Kohn–Sham theory}
With $\chi=\chi^{(1,7)}$, four real couplings $g_S,g_P,g_V,g_A$ and four real external sources $\mathsf s,\mathsf p,\mathsf v_A,\mathsf a_A$:
\begin{align*}
\mathcal L&=\bar\Psi\big(i\Gamma^{A}\partial_A-m\big)\Psi-\bar\Psi\big(\mathsf s+i\chi\,\mathsf p+\Gamma^{A}(\mathsf v_A+\chi\,\mathsf a_A)\big)\Psi \\
&\quad-\frac{g_S}2(\bar\Psi\Psi)^2-\frac{g_P}2(\bar\Psi i\chi\Psi)^2-\frac{g_V}2(\bar\Psi\Gamma^{A}\Psi)(\bar\Psi\Gamma_{A}\Psi)-\frac{g_A}2(\bar\Psi\Gamma^{A}\chi\Psi)(\bar\Psi\Gamma_{A}\chi\Psi),
\end{align*}
where $\Gamma_A=\Gamma^{(1,7)}_A$. The densities $\rho_S=\bar\Psi\Psi$, $\rho_P=\bar\Psi i\chi\Psi$, $j^A=\bar\Psi\Gamma^A\Psi$, $j_5^A=\bar\Psi\Gamma^A\chi\Psi$ (18 real fields) determine the ground state (Hohenberg–Kohn), and the Kohn–Sham orbitals $\Phi_k$ obey
\begin{align*}
H_{\rm KS}\Phi_k&=\lambda_k\Phi_k, \\
H_{\rm KS}&=\Gamma^{0}\Gamma^{i}\big(-i\partial_i+\mathcal V_i\big)+\Gamma^{0}\big(m+\mathcal S\big)+\mathcal V_0+\Gamma^{0}i\chi\,\mathcal P+\chi\,\mathcal A_0+\Gamma^{0}\Gamma^{i}\chi\,\mathcal A_i, \\
\mathcal S&=\mathsf s+g_S\rho_S+\frac{\delta E_{xc}}{\delta\rho_S},\qquad \mathcal P=\mathsf p+g_P\rho_P+\frac{\delta E_{xc}}{\delta\rho_P}, \\
\mathcal V_A&=\mathsf v_A+g_Vj_A+\frac{\delta E_{xc}}{\delta j^A},\qquad \mathcal A_A=\mathsf a_A+g_Aj_{5A}+\frac{\delta E_{xc}}{\delta j_5^A}, \\
\rho_S&=\sum_k\bar\Phi_k\Phi_k,\quad \rho_P=\sum_k\bar\Phi_ki\chi\Phi_k,\quad j^A=\sum_k\bar\Phi_k\Gamma^{A}\Phi_k,\quad j_5^A=\sum_k\bar\Phi_k\Gamma^{A}\chi\Phi_k,
\end{align*}
with $j_A=\eta^{(1,7)}_{AA}j^A$, the sums running over the occupied positive-energy orbitals (no Dirac sea), and $E_{xc}=F-T_s-E_H$, where $F$ is the universal functional, $T_s=\sum_k\int\Phi_k^\dagger H_{\rm D}\Phi_k$ and $E_H=\int\big[\frac{g_S}2\rho_S^2+\frac{g_P}2\rho_P^2+\frac{g_V}2j^Aj_A+\frac{g_A}2j_5^Aj_{5A}\big]$. The exact exchange energy of the contact couplings follows from Wick’s theorem: with the 16×16 matrix $\mathcal R=\sum_k\Phi_k\bar\Phi_k$,
\begin{equation*}
E_x=-\frac{g_S}2\mathrm{tr}(\mathcal R\mathcal R)-\frac{g_P}2\mathrm{tr}(i\chi\mathcal R\,i\chi\mathcal R)-\frac{g_V}2\mathrm{tr}(\Gamma^{A}\mathcal R\Gamma_{A}\mathcal R)-\frac{g_A}2\mathrm{tr}(\Gamma^{A}\chi\mathcal R\Gamma_{A}\chi\mathcal R).
\end{equation*}
\Needspace{8\baselineskip}\subsection*{11.4 The chiral Kohn–Sham theory}\addcontentsline{toc}{subsection}{11.4 The chiral Kohn–Sham theory}
In chiral components $\Phi_{kL}$, $\Phi_{kR}$ (as for $\Psi_L$, $\Psi_R$), with $\kappa=\sum_k\Phi_{kL}^\dagger\Phi_{kR}$ ($\rho_S=\kappa+\kappa^*$, $\rho_P=i(\kappa-\kappa^*)$), the complex Kohn–Sham mass $\mathcal M=m+\mathcal S+i\mathcal P$ and $\mathcal V^{L,R}_A=\mathcal V_A\mp\mathcal A_A$, the Kohn–Sham equation becomes
\begin{align*}
\big[\mathcal V^L_0-\tau^{O}_i\partial_i-i\tau^{O}_i\mathcal V^L_i\big]\Phi_L+\mathcal M\,\Phi_R&=\lambda\,\Phi_L, \\
\mathcal M^*\,\Phi_L+\big[\mathcal V^R_0+\tau^{O}_i\partial_i+i\tau^{O}_i\mathcal V^R_i\big]\Phi_R&=\lambda\,\Phi_R .
\end{align*}
The scalar–pseudoscalar Hartree energy $\frac{g_S-g_P}2(\kappa^2+\kappa^{*2})+(g_S+g_P)\kappa\kappa^*$ is invariant under the real chiral rotations $e^{i\vartheta\chi}=\cos\vartheta\,1_{16}-\sin\vartheta\,\Gamma^{(1,7)}_8$ ($\kappa\to e^{2i\vartheta}\kappa$) exactly when $g_P=g_S$.

\Needspace{8\baselineskip}\subsection*{11.5 Examples and verification}\addcontentsline{toc}{subsection}{11.5 Examples and verification}
\textbf{Uniform matter} (units $m=1$, no sea, $g_S=-10^5$, $g_V=10^4$, $g_P=g_A=0$): the Hartree gap equation $m^*=m+g_S\rho_S(m^*)$ and the energy density give saturation at the Fermi momentum $k_F=0.756636150702$, where the energy per particle is $E/A-m=-0.0283214279647$ (bound) and $m^*=0.348097340146$; there $E/A$ equals the Fermi energy $E_F=g_V\rho+\sqrt{k_F^2+m^{*2}}$ (Hugenholtz–Van Hove), with $\rho=k_F^7/(105\pi^4)$.

\textbf{The chiral vacuum with the Dirac sea} (all negative-energy orbitals filled, momentum cutoff $|\vec p|<\Lambda_{\rm c}=1$, $g_S=g_P=g$): for $m=0$ a non-zero mass $|\mathcal M|$ appears exactly when $g<g_c=-90\pi^4/\Lambda_{\rm c}^6$; for $g=\frac65g_c$, $|\mathcal M|=0.553693851829$ — spontaneous breaking of the chiral symmetry.

\textbf{Checks.} Sections 7–9 of 7+1.nb contain 74 checks, all True (section 9); the independent program \texttt{tools/ref\_dirac16.py} repeats the constructions with 90 checks, all True, and reproduces every printed number to at least 17 digits.

\Needspace{12\baselineskip}\section*{12. Sections 11–19: gravity, the self-coupled dirac16complex in gravitational fields, and the Standard Model}\addcontentsline{toc}{section}{12. Sections 11–19: gravity, the self-coupled dirac16complex in gravitational fields, and the Standard Model}
Sections 11–19 of 7+1.nb (and the summary of all checks, then section 20, now section 24) were appended in a third stage; sections 0–10 are unchanged, except for the explanation of the signature added on 27 October 2026 (section 2). They rebuild the ideas of the repository Dirac\_claude on the 7+1 spacetime of this project — ONE time coordinate $x_0$ and SEVEN space coordinates $x_1,\dots,x_7$ ($x_1,x_2,x_3$ ordinary space, $x_4,\dots,x_7$ auxiliary space); the 4+4 metric of Dirac\_claude is not used. Here $\Gamma_A=\Gamma^{(1,7)}_A$, $\Gamma^A=\eta_{(1,7)}^{AA}\Gamma_A$, $\eta=\eta^{(1,7)}=\mathrm{diag}(+1,-1,\dots,-1)$ — simply the squares of the octonion basis, $e_0^2=+1$, $e_i^2=-1$, $\eta_{AB}=\mathrm{Re}(e_Ae_B)$ (the tensor package xAct declares the same spacetime, as the reference notebook does, with the signature list $\{7,1,0\}$: seven space-like dimensions, one time-like dimension, none degenerate) —, $\chi=i\Gamma_0\Gamma_1\cdots\Gamma_7$, $\Psi$ is the field dirac16complex and $\bar\Psi=\Psi^\dagger\Gamma_0$; a prime is $d/dx_0$. The sections contain 100, 24 and 24 checks (gravity, Kohn–Sham, Standard Model), of which 100, 24 and 24 are True.

\Needspace{8\baselineskip}\subsection*{12.1 Curvature and the Lovelock tensors (sections 11–12)}\addcontentsline{toc}{subsection}{12.1 Curvature and the Lovelock tensors (sections 11–12)}
Christoffel symbols are written $\{{}^{\lambda}_{\mu\nu}\}$ ($\Gamma$ always denotes a Dirac matrix); the curvature follows Landau and Lifshitz:
\begin{align*}
\{{}^{\lambda}_{\mu\nu}\}&=\tfrac12g^{\lambda\rho}\big(\partial_\mu g_{\rho\nu}+\partial_\nu g_{\rho\mu}-\partial_\rho g_{\mu\nu}\big),\qquad R^{\rho}{}_{\sigma\mu\nu}=\partial_\mu\{{}^{\rho}_{\sigma\nu}\}-\partial_\nu\{{}^{\rho}_{\sigma\mu}\}+\{{}^{\rho}_{\mu\lambda}\}\{{}^{\lambda}_{\sigma\nu}\}-\{{}^{\rho}_{\nu\lambda}\}\{{}^{\lambda}_{\sigma\mu}\}, \\
R_{\sigma\nu}&=R^{\rho}{}_{\sigma\rho\nu},\qquad R=g^{\sigma\nu}R_{\sigma\nu},\qquad G^\mu{}_\nu=R^\mu{}_\nu-\tfrac12\delta^\mu_\nu R,\qquad R^{\gamma\delta}{}_{\alpha\beta}=g^{\delta\sigma}R^{\gamma}{}_{\sigma\alpha\beta}.
\end{align*}
In $n$ dimensions the most general symmetric, divergence-free tensor built from the metric and its first two derivatives is a sum of $m-1$ Lovelock tensors plus a cosmological term, $m=n/2$ for even $n$ (D. Lovelock and H. Rund): for $n=8$ there are exactly three, normalized as in the reference notebook Einstein-Lovelock-7+1-Nash.nb (P. L. Nash, GPL),
\begin{align*}
E^{(k)\mu}{}_\nu&=-\frac{1}{2^{k+1}}\,\delta^{\mu\alpha_1\beta_1\cdots\alpha_k\beta_k}_{\nu\gamma_1\delta_1\cdots\gamma_k\delta_k}\,R^{\gamma_1\delta_1}{}_{\alpha_1\beta_1}\cdots R^{\gamma_k\delta_k}{}_{\alpha_k\beta_k},\qquad L_k=\frac{1}{2^k}\,\delta^{\alpha_1\beta_1\cdots\alpha_k\beta_k}_{\gamma_1\delta_1\cdots\gamma_k\delta_k}\,R^{\gamma_1\delta_1}{}_{\alpha_1\beta_1}\cdots R^{\gamma_k\delta_k}{}_{\alpha_k\beta_k}, \\
\delta^{\mu_1\cdots\mu_n}_{\nu_1\cdots\nu_n}&=\det\big[\delta^{\mu_i}_{\nu_j}\big],\qquad E^{(1)}=G,\qquad E^{(k)\mu}{}_\mu=-\frac{8-2k}{2}L_k,\qquad E^{(4)}\equiv0 .
\end{align*}
The field equations, in the form of the reference notebook (pure numbers $w_k$, $\Lambda$ and an inverse length $H$) and with matter ($\kappa=8\pi G_8$):
\begin{gather*}
0=-\Lambda\,\delta^\mu_\nu+\sum_{k=1}^{3}H^{-2k}w_k\,E^{(k)\mu}{}_\nu,\qquad \sum_{k=0}^{3}\alpha_kE^{(k)\mu}{}_\nu=\kappa T^\mu{}_\nu\ \ (\alpha_0=2\bar\Lambda,\ \alpha_1=1),\\ \alpha_2=\frac{w_2}{w_1H^2},\ \ \alpha_3=\frac{w_3}{w_1H^4},\ \ \bar\Lambda=\frac{\Lambda H^2}{w_1}.
\end{gather*}
The notebook computes the Lovelock tensors twice: with the tensor packages xAct (xPerm 1.2.3, xTensor 1.2.0, xCoba 0.8.6, xTras 1.4.2) exactly as the reference notebook does — abstract indices, the generalized delta \texttt{kδ[lower, upper] = Det[Outer[delta, lower, upper]]}, and the components of a given metric from \texttt{CD = LC[g]} — and by an independent sum over sets of index pairs; the two agree exactly for the cosmological metric and a gravitational wave. The reference notebook’s section "FAILED Model metric g71; DO NOT USE!" is not used. For the cosmological metric $ds^2=dx_0^2-a^2(dx_1^2+dx_2^2+dx_3^2)-b^2(dx_4^2+\cdots+dx_7^2)$ with $H_a=a'/a$, $H_b=b'/b$:
\begin{align*}
E^{(1)0}{}_0&=3H_a^2+12H_aH_b+6H_b^2,\qquad E^{(2)0}{}_0=-12H_b\big(4H_a^3+18H_a^2H_b+12H_aH_b^2+H_b^3\big), \\
E^{(3)0}{}_0&=360\,H_a^2H_b^3\big(4H_a+3H_b\big).
\end{align*}
Exponential vacuum solutions ($a=e^{Hx_0}$, $b=e^{q_HHx_0}$) exist for a one-parameter family of couplings; for $q_H=-1$ (eternal inflation of ordinary space with eternal deflation of the auxiliary space) $w_2=\frac{\Lambda}{4} - \frac{w_{1}}{2}$, $w_3=- \frac{\Lambda}{36} + \frac{w_{1}}{24}$.

\Needspace{8\baselineskip}\subsection*{12.2 Generalized frame matrices, the canonical spin connection, and the covariant Dirac equation (section 13)}\addcontentsline{toc}{subsection}{12.2 Generalized frame matrices, the canonical spin connection, and the covariant Dirac equation (section 13)}
With a frame field $g_{\mu\nu}=\eta_{AB}e^A{}_\mu e^B{}_\nu$ and its inverse $e_A{}^\mu$:
\begin{align*}
\Gamma_\mu(X)&=e^A{}_\mu\Gamma_A,\qquad \Gamma^\mu(X)=e_A{}^\mu\Gamma^A,\qquad \Gamma_\mu\Gamma_\nu+\Gamma_\nu\Gamma_\mu=2g_{\mu\nu}1_{16}, \\
\omega_\mu{}^A{}_B&=e^A{}_\nu\big(\partial_\mu e_B{}^\nu+\{{}^{\nu}_{\mu\lambda}\}e_B{}^\lambda\big),\qquad \Omega_\mu=\tfrac14\eta_{AC}\,\omega_\mu{}^C{}_B\,\Gamma^A\Gamma^B,\qquad D_\mu=\partial_\mu+\Omega_\mu, \\
\nabla_\mu\Gamma_\nu&=\partial_\mu\Gamma_\nu-\{{}^{\lambda}_{\mu\nu}\}\Gamma_\lambda+[\Omega_\mu,\Gamma_\nu]=0,\qquad [D_\mu,D_\nu]=\tfrac14R_{AB\mu\nu}\Gamma^A\Gamma^B,\qquad (\Gamma^\mu D_\mu)^2=g^{\mu\nu}\nabla_\mu D_\nu-\tfrac R4 .
\end{align*}
The covariant Dirac equation $\big(i\Gamma^\mu(X)D_\mu-m\big)\Psi=0$ follows from $\mathcal L=\sqrt{|g|}\,\bar\Psi(i\Gamma^\mu D_\mu-m)\Psi$, and $\bar\Psi\Gamma^\mu\Psi$ is conserved. For the cosmological metric the spin connection is not zero:
\begin{equation*}
\Omega_0=0,\quad \Omega_i=\tfrac{a'}{2}\Gamma_i\Gamma_0\ (i=1,2,3),\quad \Omega_k=\tfrac{b'}{2}\Gamma_k\Gamma_0\ (k=4,\dots,7),\qquad \Gamma^\mu\Omega_\mu=\big(\tfrac32H_a+2H_b\big)\Gamma_0 .
\end{equation*}
The notebook also derives the connection from Cartan’s structure equation $de^A+\omega^A{}_B\wedge e^B=0$ for a static metric with a radial coordinate, and couples an abelian gauge field, $D_\mu=\partial_\mu+\Omega_\mu-iqA_\mu$.

\Needspace{8\baselineskip}\subsection*{12.3 The three epochs (section 14)}\addcontentsline{toc}{subsection}{12.3 The three epochs (section 14)}
\textbf{Primordial} (near $x_0=0$): for $a=x_0^p$, $b=x_0^q$ the highest Lovelock order dominates. Einstein: $3p+4q=1$, $3p^2+4q^2=1$, $p=(1\pm2\sqrt2)/7$, $q=(2\mp3\sqrt2)/14$ — on the branch $p=0.547$, $q=-0.160$ ordinary space expands while the auxiliary space contracts. Gauss–Bonnet: $3p+4q=3$, $-203p^3+81p^2+351p+27=0$ ($p=-1.08350$, $-0.0786311$, $1.56115$). Cubic: $(p,q)=(-15/7,20/7)$. The anisotropic expansion creates pairs of dirac16complex. \textbf{Modern}: with $b=b_0$ the field equations reduce to the Friedmann equations of ordinary space, $3H_a^2=\kappa\varepsilon+\bar\Lambda$, $2H_a'+3H_a^2=-\kappa p_a+\bar\Lambda$, with $G_4=G_8/V_\perp$. \textbf{Far future}: in de Sitter space pairs are created with the Fermi–Dirac occupation at the Gibbons–Hawking temperature $H/2\pi$:
\begin{equation*}
|\beta|^2=\frac{1}{e^{2\pi M/H}+1}.
\end{equation*}
\Needspace{8\baselineskip}\subsection*{12.4 The nine-channel self-coupling and Kohn–Sham theory in gravitational fields (sections 15–16)}\addcontentsline{toc}{subsection}{12.4 The nine-channel self-coupling and Kohn–Sham theory in gravitational fields (sections 15–16)}
With $\Theta^K=i^{\sigma_r}\Gamma^{A_1}\cdots\Gamma^{A_r}$ for the 256 increasing multi-indices $K$ ($\sigma_r=\frac{r(r-1)}2\bmod2$, $\Theta_K=\eta_{A_1A_1}\cdots\eta_{A_rA_r}\Theta^K$) the most general contact self-coupling has one coupling per rank,
\begin{equation*}
\mathcal L_{\rm self}=-\frac12\sum_{r=0}^{8}g_r\sum_{|K|=r}(\bar\Psi\Theta_K\Psi)(\bar\Psi\Theta^K\Psi),\qquad \big(i\Gamma^\mu D_\mu-m\big)\Psi=\sum_rg_r\sum_{|K|=r}(\bar\Psi\Theta_K\Psi)\,\Theta^K\Psi,
\end{equation*}
and contains the four couplings of section 11.3 ($g_0=g_S$, $g_1=g_V$, $g_7=-g_A$, $g_8=-g_P$). The Fierz matrix $F_{rs}=\frac1{16}\sum_l(-1)^{rs-l}\binom sl\binom{8-s}{r-l}$ ($F^2=1$) turns exchange into direct terms, so Hartree–Fock is Hartree with $G=(1-F^{\mathsf T})g$. Homogeneous Kohn–Sham matter in the cosmological spacetime has the effective mass $M=m+G_0\rho_S$, the energy density $\varepsilon$ and the pressures $p_a$, $p_b$ (integrals over the comoving Fermi ball), and
\begin{equation*}
\varepsilon'+3H_a(\varepsilon+p_a)+4H_b(\varepsilon+p_b)=0 .
\end{equation*}
Near $x_0=0$ the effective mass falls toward zero ($M=0.003824$ at $x_0=10^{-4}$), in the modern epoch the matter is pressureless in ordinary space, and in a de Sitter future $1-M$ decays like the density; Einstein and Einstein–Lovelock gravity driven by this matter preserve the constraint equation to a relative $2.4\times10^{-14}$.

\Needspace{8\baselineskip}\subsection*{12.5 The Standard Model on the world brane (section 17)}\addcontentsline{toc}{subsection}{12.5 The Standard Model on the world brane (section 17)}
The auxiliary coordinates form a torus ($x_k\simeq x_k+L_c$); the Standard Model lives on the brane $W=\lbrace x_4=\dots=x_7=0\rbrace$ with the induced metric $h_{\mu\nu}=g_{\mu\nu}(x,y=0)$, and the total action is
\begin{equation*}
S=S_{\rm EL}[g]+S_\Psi[\Psi,g]+\int_Wd^4x\sqrt{|h|}\,\big(\mathcal L_{\rm SM}+\mathcal L_{\rm portal}\big)+S_{\rm stab},
\end{equation*}
\begin{equation*}
S_{\rm EL}=-\frac1{2\kappa}\int d^8X\sqrt{|g|}\sum_{k=0}^{3}\alpha_kL_k,\qquad S_\Psi=\int d^8X\sqrt{|g|}\,\big[\bar\Psi(i\Gamma^\mu D_\mu-m)\Psi+\mathcal L_{\rm self}\big],
\end{equation*}
($S_{\rm stab}$ is a stress that keeps the auxiliary scale factor $b$ constant in the modern epoch), with the complete Standard Model ($SU(3)\times SU(2)\times U(1)$, three generations with right-handed neutrinos, one Higgs doublet, Yukawa couplings) in the gravitational field of $h$. All gauge anomalies cancel in every generation; with the data of the Particle Data Group the electroweak scale is $v=(\sqrt2G_F)^{-1/2}=246.219640$ GeV and the CKM Jarlskog invariant $J=3.117\times10^{-5}$; the brane Friedmann equation $3H_a^2=8\pi G_4\rho_4+\Lambda_4$ with the Planck 2018 parameters gives the age 13.791 Gyr.

\Needspace{8\baselineskip}\subsection*{12.6 dirac16complex at the brane, portals, and neutrino flavour (sections 18–19)}\addcontentsline{toc}{subsection}{12.6 dirac16complex at the brane, portals, and neutrino flavour (sections 18–19)}
With the Dirac representation $\gamma_\mu$ ($\gamma_\mu\gamma_\nu+\gamma_\nu\gamma_\mu=2\eta_{\mu\nu}1_4$), $\gamma_5=i\gamma_0\gamma_1\gamma_2\gamma_3$ and $\rho_4=i\gamma_0$, $\rho_5=-\gamma_1$, $\rho_6=-\gamma_2$, $\rho_7=-\gamma_3$, the matrices
\begin{equation*}
R_\mu=\gamma_\mu\otimes1_4\ (\mu=0,\dots,3),\qquad R_k=\gamma_5\otimes\rho_k\ (k=4,\dots,7),\qquad \Gamma_AU=UR_A,\quad U^\dagger U=1_{16},
\end{equation*}
show that dirac16complex at the brane is four ordinary Dirac spinors, $\Psi=U\sum_a\psi_a\otimes e_a$, with $\bar\Psi\Gamma^\mu\Psi'=\sum_a\bar\psi_a\gamma^\mu\psi'_a$; a mode with auxiliary momentum $\vec q$ has the Kaluza–Klein mass $\sqrt{m^2+|\vec q|^2/b_0^2}$. Three portals (Yukawa, vector, Higgs) couple it to the Standard Model at the brane; the Yukawa portal makes the light neutrino mixing non-unitary, $1-NN^\dagger=\mu\mu^\dagger/m^2$. Neutrino flavour conversion in vacuum,
\begin{equation*}
P(\alpha\to\beta)=\delta_{\alpha\beta}-4\sum_{i>j}\mathrm{Re}\,X_{ij}\sin^2\frac{\Delta m^2_{ij}L}{4E}+2\sum_{i>j}\mathrm{Im}\,X_{ij}\sin\frac{\Delta m^2_{ij}L}{2E},\qquad X_{ij}=U^*_{\alpha i}U_{\beta i}U_{\alpha j}U^*_{\beta j},
\end{equation*}
in matter (potential $\sqrt2G_Fn_e$), in the Sun and in supernovae is worked out with NuFit-6.0 data; for DUNE ($L=1300$ km, $E=2.5$ GeV) $P(\nu_\mu\to\nu_e)=0.048837$ in vacuum and 0.072327 in the Earth’s crust.

\textbf{Checks.} Sections 11–19 of 7+1.nb contain 148 checks, all True; the independent programs \texttt{tools/ref\_gravity.py}, \texttt{tools/ref\_ks\_gravity.py} and \texttt{tools/ref\_sm.py} repeat the constructions with 101, 29 and 33 checks, all True, and reproduce the printed numbers.

\Needspace{12\baselineskip}\section*{13. Sections 20–23: the coupled field equations and their solutions in the three gravity regimes}\addcontentsline{toc}{section}{13. Sections 20–23: the coupled field equations and their solutions in the three gravity regimes}
Sections 20–23 of 7+1.nb were appended in a fourth stage (27 October 2026); sections 0–19 are unchanged, and the summary of all checks, formerly section 20, is now section 24. The coordinates are those of the whole project: $x_0$ time (the octonion unit), $x_1,x_2,x_3$ ordinary space, $x_4,\dots,x_7$ auxiliary space; $\eta=\mathrm{diag}(+1,-1,\dots,-1)$ is simply the squares of the octonion basis, and xAct declares the spacetime with the signature list $\{7,1,0\}$ (seven space-like, one time-like dimension). A prime is $d/dx_0$; $\kappa=8\pi G_8=M_8^{-6}$. The sections contain 48 checks, of which 48 are True.

\Needspace{8\baselineskip}\subsection*{13.1 Every interaction in equation form (section 20)}\addcontentsline{toc}{subsection}{13.1 Every interaction in equation form (section 20)}
With the Lovelock densities $L_k=2^{-k}\delta^{\alpha_1\beta_1\cdots\alpha_k\beta_k}_{\gamma_1\delta_1\cdots\gamma_k\delta_k}R^{\gamma_1\delta_1}{}_{\alpha_1\beta_1}\cdots R^{\gamma_k\delta_k}{}_{\alpha_k\beta_k}$ ($\delta^{\cdots}_{\cdots}$ the generalized Kronecker delta), the real 16×16 Dirac matrices $\Gamma^A$, the frame matrices $\Gamma^\mu=e_A{}^\mu\Gamma^A$, the spin connection $\Omega_\mu=\frac14\eta_{AC}\omega_\mu{}^C{}_B\Gamma^A\Gamma^B$, the Clifford products $\Theta^K=i^{\sigma_r}\Gamma^{A_1}\cdots\Gamma^{A_r}$ ($\sigma_r=\frac{r(r-1)}2\bmod2$, $\eta_K=\eta_{A_1A_1}\cdots\eta_{A_rA_r}$), the world brane $W=\lbrace x_4=\dots=x_7=0\rbrace$ with induced metric $h$, and $b^8=\det(-g_{kl})$ ($k,l=4,\dots,7$), the action is
\begin{align*}
S&=-\frac1{2\kappa}\int d^8X\sqrt{|g|}\big(2\Lambda+R+\alpha_2L_2+\alpha_3L_3\big) \\
&\quad+\int d^8X\sqrt{|g|}\Big[\bar\psi\,i\Gamma^\mu(\partial_\mu+\Omega_\mu)\psi-m\bar\psi\psi \\
&\qquad\qquad-\frac12\sum_{r=0}^8g_r\sum_{|K|=r}\eta_K(\bar\psi\Theta^K\psi)^2-U(b)\Big] \\
&\quad+\int_Wd^4x\sqrt{|h|}\Big[-f_W+\mathcal L_{\rm SM}-\sum_{i,s}y_{is}\bar L_i\tilde\Phi P_R\psi_s+\text{h.c.} \\
&\qquad\qquad-G^f_P(\bar\psi\Gamma^\mu\psi)(\bar f\gamma_\mu f)-\lambda_P\Phi^\dagger\Phi\,\bar\psi\psi\Big], \\
L_2&=R^2-4R_{\mu\nu}R^{\mu\nu}+R_{\mu\nu\rho\sigma}R^{\mu\nu\rho\sigma},\qquad U(b)=U_0+\frac{K_U}2\ln^2\frac b{b_0}, \\
L_3&=R^3-12RR_{\mu\nu}R^{\mu\nu}+16R^\mu{}_\nu R^\nu{}_\rho R^\rho{}_\mu+24R^{\mu\rho}R^{\nu\sigma}R_{\mu\nu\rho\sigma}+3RR_{\mu\nu\rho\sigma}R^{\mu\nu\rho\sigma} \\
&\quad-24R^{\mu\nu}R_\mu{}^{\rho\sigma\tau}R_{\nu\rho\sigma\tau}+4R^{\mu\nu}{}_{\rho\sigma}R^{\rho\sigma}{}_{\lambda\tau}R^{\lambda\tau}{}_{\mu\nu}-8R_\mu{}^\rho{}_\nu{}^\sigma R_\rho{}^\lambda{}_\sigma{}^\tau R_\lambda{}^\mu{}_\tau{}^\nu,
\end{align*}
where $\mathcal L_{\rm SM}$ is the Standard Model Lagrangian (gauge fields $G^a_\mu,W^I_\mu,B_\mu$, Higgs doublet $\Phi$, three generations, Yukawa couplings) in the gravitational field of $h$, $y_{is}$, $G^f_P$, $\lambda_P$ are the Yukawa, vector and Higgs portal couplings, $L_i$ the lepton doublets and $\psi_s$ the brane-right-handed parts of dirac16complex at the brane. Varying it gives $\sum_k\alpha_kE^{(k)\mu}{}_\nu=\kappa T^\mu{}_\nu$ ($\alpha_0=2\Lambda$, $\alpha_1=1$, $E^{(k)\mu}{}_\nu=-2^{-k-1}\delta^{\mu\alpha_1\cdots}_{\nu\gamma_1\cdots}R\cdots R$), the nonlinear Dirac equation of $\psi$ with the portal sources at the brane, and the Standard Model equations with portal terms. For the homogeneous metric every source is $T^\mu{}_\nu=\mathrm{diag}(\varepsilon,-p_a\times3,-p_b\times4)$, and for a comoving energy $W(a,b)$
\begin{align*}
\varepsilon&=\frac W{a^3b^4},\quad p_a=-\frac{\partial_aW}{3a^2b^4},\quad p_b=-\frac{\partial_bW}{4a^3b^3}, \\
\beta''+3H_a\beta'+4\beta'^2&=\kappa\Big(\frac\varepsilon6-\frac{p_a}2+\frac{p_b}3\Big)+\frac\Lambda3\quad(\beta=\ln\tfrac b{b_0},\ \alpha_2=\alpha_3=0),
\end{align*}
so the radion $\beta$ has the mass $m_r^2=\kappa K_U/12$, and with $b$ static the brane obeys $3H_a^2=8\pi G_4\rho_4+\Lambda_4$ with $8\pi G_4=\kappa/V_\perp$ ($V_\perp=b^4L_c^4$ the volume of the auxiliary torus) and $\Lambda_4=\Lambda+\kappa U_0+\kappa f_W/V_\perp$. The canonical radion $\phi=\sqrt{12}M_{\rm Pl}\beta$ couples to the brane through $(\phi/(\sqrt3M_{\rm Pl}))T^\mu{}_\mu$.

\Needspace{8\baselineskip}\subsection*{13.2 Regime I: the primordial Lovelock regime (section 21)}\addcontentsline{toc}{subsection}{13.2 Regime I: the primordial Lovelock regime (section 21)}
In units $\ell_2=\sqrt{\alpha_2}=1$ with $\xi=\alpha_3/\alpha_2^2=10^{-6}$ and $u=H_b/H_a$, $X=H_a^2$ ($H_a=a'/a$, $H_b=b'/b$), the vacuum constraint is quadratic in $X$:
\begin{align*}
0&=3(2u^2+4u+1)-12u(u^3+12u^2+18u+4)X+360\xi u^3(3u+4)X^2, \\
\tau(H_a)&=\int_{H_a}^\infty\frac{dH}{F(H)},\qquad F=-\frac{dH_a}{dx_0}>0,
\end{align*}
an exact solution as an explicit curve plus one quadrature ($\tau$ the time since the singularity). Its limits: the cubic regime $H_a=1/\tau+A_0+A_1\tau$, $H_b=h_0+h_1\tau+h_2\tau^2$ with $h_0=-\sqrt{\alpha_2/(30\alpha_3)}=-182.5742$; the Gauss–Bonnet Kasner regime $a\propto\tau^{p_2}$, $b\propto\tau^{q_2}$ ($p_2=3/(3+4u_2)$, $u_2^3+12u_2^2+18u_2+4=0$); the Einstein Kasner regime $p_1=(1+2\sqrt2)/7=0.5469182$, $q_1=(2-3\sqrt2)/14=-0.1601886$, with $\tau=p_1/H_a-e_1H_a+s_1$, $e_1=\alpha_2(1+9\sqrt2)/7$. Inside their regimes these formulas are accurate to $6.7\times10^{-6}$ (cubic), 0.00085 (Gauss–Bonnet) and 0.00032 (Einstein). With a dirac16complex gas, $p_a=w_a\varepsilon$, the universe near $\tau=0$ is matter-dominated: $a\propto\tau^{1/(1+w_a)}$ ($\tau^{3/4}$ for the gas; measured 0.7500003) and $H_b\to h_*=-\sqrt{\alpha_2/(54\alpha_3)}$. Once the gas dominates the Einstein regime, the auxiliary space re-expands, $b\propto\tau^{2/5}$ (measured $\tau H_b=0.399995$): a stabilizing potential is needed.

\Needspace{8\baselineskip}\subsection*{13.3 Regime II: the modern regime (section 22)}\addcontentsline{toc}{subsection}{13.3 Regime II: the modern regime (section 22)}
The potential catches the auxiliary space when $H\approx m_r$; the radion then oscillates as
\begin{equation*}
\beta\approx\frac{A\sin(m_r\tau+\vartheta)}{m_r(\tau-\tau_E)},\qquad H_a+2H_b\approx\frac2{3(\tau-\tau_E)},
\end{equation*}
(fitted: $A=0.46933$, largest deviation 0.013 of the amplitude) and carries 0.991 of the energy. Into brane fields a 10 GeV radion ($M_8=10$ TeV) decays through $\Gamma(\phi\to f\bar f)=N_cm_f^2m_r(1-4m_f^2/m_r^2)^{3/2}/(24\pi M_{\rm Pl}^2)$ and $\Gamma(\phi\to gg)=b_{\rm QCD}^2\alpha_s^2m_r^3/(96\pi^3M_{\rm Pl}^2)$, which alone would reheat the brane only to $1.01\times10^{-9}$ GeV; but its $10^{13}$ decay channels of gravitational strength into Kaluza–Klein pairs of dirac16complex ($\Gamma_n=3\mu_n^2m_r(1-4\mu_n^2/m_r^2)^{3/2}/(8\pi M_{\rm Pl}^2)$, $\mu_n=2\pi|n|/V_\perp^{1/4}$) make it decay into the BULK within 0.11 s. The minimal model therefore has a moduli problem and needs a later reheating of the brane. After it, the brane plasma ($g_*(T)$ from ideal gases of all Standard Model particles) gives $t(1\,\mathrm{MeV})=0.74319$ s, and the late universe ($H_a=H_0\sqrt{\Omega_ra^{-4}+\Omega_ma^{-3}+\Omega_\Lambda}$, Planck 2018) has the composite analytic solution
\begin{equation*}
t(a)\approx\frac{2}{3H_0\Omega_m^2}\Big[(\Omega_ma-2\Omega_r)\sqrt{\Omega_r+\Omega_ma}+2\Omega_r^{3/2}\Big]+\frac{2\,\mathrm{arcsinh}\big(\sqrt{\Omega_\Lambda/\Omega_m}\,a^{3/2}\big)}{3H_0\sqrt{\Omega_\Lambda}}-\frac{2a^{3/2}}{3H_0\sqrt{\Omega_m}},
\end{equation*}
accurate to 0.00022 at all times; the age is 13.791 Gyr. The radion sits at $\beta_{\rm eq}=1.7\times10^{-86}$ today, so $G_4$ does not change.

\Needspace{8\baselineskip}\subsection*{13.4 Regime III: the far future (section 23)}\addcontentsline{toc}{subsection}{13.4 Regime III: the far future (section 23)}
De Sitter expansion with $H_\infty=H_0\sqrt{\Omega_\Lambda}=1.808\times10^{-18}$ s$^{-1}$, exactly $H_a=H_\infty\coth(\tfrac32H_\infty t)$ for matter and $\Lambda$; the radion obeys $\beta''+\gamma\beta'+m_{\rm eff}^2(\beta-\beta_{\rm eq})=0$ with
\begin{align*}
\frac{\gamma}{3H}&=\frac{1-16\alpha_2H^2+48\alpha_2^2H^4-48\alpha_3H^4+512\alpha_2^3H^6+192\alpha_2\alpha_3H^6-1024\alpha_2^4H^8}{(1-4\alpha_2H^2+32\alpha_2^2H^4)^2}, \\
\frac{m_{\rm eff}^2}{m_r^2}&=\frac1{1-4\alpha_2H^2+32\alpha_2^2H^4},
\end{align*}
so the auxiliary space stays frozen; without the potential the far future would be eight-dimensional de Sitter, $21H^2-420\alpha_2H^4+2520\alpha_3H^6=\Lambda$, with an inflating auxiliary space. \textbf{Checks.} Sections 20–23 contain 48 checks, all True; \texttt{tools/ref\_coupled.py} repeats them with 49 checks, all True.

\Needspace{12\baselineskip}\section*{14. Files in this repository}\addcontentsline{toc}{section}{14. Files in this repository}
{\small
\begin{longtable}{>{\raggedright\arraybackslash}p{\dimexpr 0.4787\linewidth-2\tabcolsep\relax}>{\raggedright\arraybackslash}p{\dimexpr 0.5213\linewidth-2\tabcolsep\relax}}
\toprule
\textbf{path} & \textbf{content} \\
\midrule\endfirsthead
\toprule
\textbf{path} & \textbf{content} \\
\midrule\endhead
\endfoot
\bottomrule\endlastfoot
\texttt{7+1.nb} & the new notebook, evaluated, with all outputs \\
\texttt{4+4.nb} & the 4+4 notebook (split octonions, $T_A$), the starting point; unchanged \\
\texttt{NASH\_Mathematica\_Notebooks/} & earlier notebooks of the project; unchanged \\
\texttt{tools/make\_7p1\_nb.py}, \texttt{tools/nbwrite.py} & generator of 7+1.nb \\
\texttt{tools/run\_eval.wl} & evaluates 7+1.nb in place in a fresh kernel of Wolfram 15.0.1 and saves it \\
\texttt{tools/watch\_requests.wl}, \texttt{tools/git\_push.wl} & runs Wolfram 15.0.1 jobs requested in build/request.txt (evaluate, push) without keyboard input; pushes with the local Git \\
\texttt{tools/eval\_cells\_log.wl} & development check: evaluates a generated notebook cell by cell and logs messages and checks \\
\texttt{tools/ref\_dirac16.py} & independent Python checks of sections 7–9 (dirac16complex, Kohn–Sham, chiral Kohn–Sham) \\
\texttt{tools/ref44.py}, \texttt{tools/ref71.py}, \texttt{tools/extra71.py}, \texttt{tools/octalg.py} & independent Python (sympy) implementation and checks \\
\texttt{tools/classify\_clifford.py} & numerical classification of the nine signatures via commutants of realified representations \\
\texttt{tools/cayley\_dickson\_formulas.py}, \texttt{tools/cayley\_dickson\_match.py} & tests of Cayley–Dickson formulas against the tables \\
\texttt{tools/nbparse.py}, \texttt{tools/verify\_7p1\_nb.py} & parser for .nb files; cross-check of the evaluated 7+1.nb \\
\texttt{tools/docgen.py}, \texttt{tools/doc\_*.py}, \texttt{tools/make\_docs.py} & generators of the documentation (Markdown, LaTeX, PDF) \\
\texttt{tools/nb\_part\_gravity.py}, \texttt{tools/nb\_part\_ks\_gravity.py}, \texttt{tools/nb\_part\_sm.py} & generator parts of sections 11–19 of 7+1.nb (executed by \texttt{make\_7p1\_nb.py}) \\
\texttt{tools/lovelock71.py} & independent curvature and Lovelock engine (pair-set algorithm, sympy) \\
\texttt{tools/ref\_gravity.py}, \texttt{tools/ref\_ks\_gravity.py}, \texttt{tools/ref\_sm.py} & independent Python checks of sections 11–14, 15–16 and 17–19 \\
\texttt{tools/nb\_part\_coupled.py} & generator part of sections 20–23 of 7+1.nb (executed by \texttt{make\_7p1\_nb.py}) \\
\texttt{tools/ref\_coupled.py} & independent Python checks of sections 20–23 \\*
\end{longtable}}
\Needspace{12\baselineskip}\section*{15. References}\addcontentsline{toc}{section}{15. References}
\begin{itemize}
\item J. C. Baez, \emph{The Octonions}, Bull. Amer. Math. Soc. 39 (2002) 145–205, arXiv:math/0105155 — section 2.2 gives the Cayley–Dickson product $(a,b)(c,d)=(ac-d\bar b,\ \bar a d+cb)$ used in section 6.5.
\item H. B. Lawson and M.-L. Michelsohn, \emph{Spin Geometry}, Princeton University Press (1989) — classification of real Clifford algebras and the periodicity $Cl(p+1,q+1)\cong Cl(p,q)\otimes\mathbb R(2)$.
\item F. Reese Harvey, \emph{Spinors and Calibrations}, Academic Press (1990) — octonions, triality and spinors.
\item P. A. M. Dirac, Proc. R. Soc. Lond. A 117 (1928) 610 — the Dirac equation $(i\gamma^\mu\partial_\mu-m)\psi=0$, written in section 11 with the eight real $\Gamma^{(1,7)}_A$.
\item P. Hohenberg and W. Kohn, Phys. Rev. 136 (1964) B864; W. Kohn and L. J. Sham, Phys. Rev. 140 (1965) A1133 — the density determines the ground state; non-interacting Kohn–Sham orbitals reproduce the density (section 11.3).
\item A. K. Rajagopal and J. Callaway, Phys. Rev. B 7 (1973) 1912; R. M. Dreizler and E. K. U. Gross, \emph{Density Functional Theory}, Springer (1990) — relativistic density-functional theory, the no-sea approximation.
\item G. C. Wick, Phys. Rev. 80 (1950) 268; M. Fierz, Z. Phys. 104 (1937) 553 — the exchange energy of contact couplings.
\item J. D. Walecka, Ann. Phys. 83 (1974) 491; N. M. Hugenholtz and L. Van Hove, Physica 24 (1958) 363 — saturation of uniform matter, and $E/A=E_F$ at saturation (section 11.5).
\item Y. Nambu and G. Jona-Lasinio, Phys. Rev. 122 (1961) 345 — spontaneous chiral symmetry breaking by contact couplings with a cutoff (section 11.5).
\item D. Lovelock, J. Math. Phys. 12 (1971) 498; D. Lovelock and H. Rund, \emph{Tensors, Differential Forms, and Variational Principles}, Dover (1989), Remark 2 after Eq. (4.38) — the Lovelock tensors (section 12.1).
\item J. M. Martín-García et al., xAct: efficient tensor computer algebra for the Wolfram Language (xPerm, xTensor, xCoba), GPL, www.xact.es; T. Nutma, Comput. Phys. Commun. 185 (2014) 1719 (xTras) — the abstract-index computation of section 12.1.
\item L. D. Landau and E. M. Lifshitz, \emph{The Classical Theory of Fields}, Pergamon — curvature conventions; E. Kasner, Am. J. Math. 43 (1921) 217 — power-law solutions; G. W. Gibbons and S. W. Hawking, Phys. Rev. D 15 (1977) 2738 — the de Sitter temperature $H/2\pi$.
\item S. Navas et al. (Particle Data Group), Phys. Rev. D 110 (2024) 030001 and the 2025 update; I. Esteban et al., JHEP 12 (2024) 216 (NuFit-6.0); Planck Collaboration, Astron. Astrophys. 641 (2020) A6 — the data of section 12.5–12.6.
\item N. Arkani-Hamed, S. Dimopoulos, G. Dvali, Phys. Lett. B 429 (1998) 263; L. Wolfenstein, Phys. Rev. D 17 (1978) 2369; S. P. Mikheyev and A. Yu. Smirnov, Sov. J. Nucl. Phys. 42 (1985) 913 — extra dimensions, the MSW effect (section 12.5–12.6).
\item D. Lovelock, J. Math. Phys. 12 (1971) 498; F. Müller-Hoissen, Phys. Lett. B 163 (1985) 106 — Lovelock gravity, the cubic density (section 13); N. Deruelle and L. Fariña-Busto, Phys. Rev. D 41 (1990) 3696 — Lovelock cosmology.
\item G. D. Coughlan, W. Fischler, E. W. Kolb, S. Raby and G. G. Ross, Phys. Lett. B 131 (1983) 59 — the cosmological moduli problem (section 13.3); N. Arkani-Hamed, S. Dimopoulos and G. Dvali, Phys. Rev. D 59 (1999) 086004 — the normalcy temperature.
\end{itemize}
\end{document}
